% Options for packages loaded elsewhere
\PassOptionsToPackage{unicode}{hyperref}
\PassOptionsToPackage{hyphens}{url}
\documentclass[
]{article}
\usepackage{xcolor}
\usepackage[margin=1in]{geometry}
\usepackage{amsmath,amssymb}
\setcounter{secnumdepth}{-\maxdimen} % remove section numbering
\usepackage{iftex}
\ifPDFTeX
  \usepackage[T1]{fontenc}
  \usepackage[utf8]{inputenc}
  \usepackage{textcomp} % provide euro and other symbols
\else % if luatex or xetex
  \usepackage{unicode-math} % this also loads fontspec
  \defaultfontfeatures{Scale=MatchLowercase}
  \defaultfontfeatures[\rmfamily]{Ligatures=TeX,Scale=1}
\fi
\usepackage{lmodern}
\ifPDFTeX\else
  % xetex/luatex font selection
\fi
% Use upquote if available, for straight quotes in verbatim environments
\IfFileExists{upquote.sty}{\usepackage{upquote}}{}
\IfFileExists{microtype.sty}{% use microtype if available
  \usepackage[]{microtype}
  \UseMicrotypeSet[protrusion]{basicmath} % disable protrusion for tt fonts
}{}
\makeatletter
\@ifundefined{KOMAClassName}{% if non-KOMA class
  \IfFileExists{parskip.sty}{%
    \usepackage{parskip}
  }{% else
    \setlength{\parindent}{0pt}
    \setlength{\parskip}{6pt plus 2pt minus 1pt}}
}{% if KOMA class
  \KOMAoptions{parskip=half}}
\makeatother
\usepackage{color}
\usepackage{fancyvrb}
\newcommand{\VerbBar}{|}
\newcommand{\VERB}{\Verb[commandchars=\\\{\}]}
\DefineVerbatimEnvironment{Highlighting}{Verbatim}{commandchars=\\\{\}}
% Add ',fontsize=\small' for more characters per line
\usepackage{framed}
\definecolor{shadecolor}{RGB}{248,248,248}
\newenvironment{Shaded}{\begin{snugshade}}{\end{snugshade}}
\newcommand{\AlertTok}[1]{\textcolor[rgb]{0.94,0.16,0.16}{#1}}
\newcommand{\AnnotationTok}[1]{\textcolor[rgb]{0.56,0.35,0.01}{\textbf{\textit{#1}}}}
\newcommand{\AttributeTok}[1]{\textcolor[rgb]{0.13,0.29,0.53}{#1}}
\newcommand{\BaseNTok}[1]{\textcolor[rgb]{0.00,0.00,0.81}{#1}}
\newcommand{\BuiltInTok}[1]{#1}
\newcommand{\CharTok}[1]{\textcolor[rgb]{0.31,0.60,0.02}{#1}}
\newcommand{\CommentTok}[1]{\textcolor[rgb]{0.56,0.35,0.01}{\textit{#1}}}
\newcommand{\CommentVarTok}[1]{\textcolor[rgb]{0.56,0.35,0.01}{\textbf{\textit{#1}}}}
\newcommand{\ConstantTok}[1]{\textcolor[rgb]{0.56,0.35,0.01}{#1}}
\newcommand{\ControlFlowTok}[1]{\textcolor[rgb]{0.13,0.29,0.53}{\textbf{#1}}}
\newcommand{\DataTypeTok}[1]{\textcolor[rgb]{0.13,0.29,0.53}{#1}}
\newcommand{\DecValTok}[1]{\textcolor[rgb]{0.00,0.00,0.81}{#1}}
\newcommand{\DocumentationTok}[1]{\textcolor[rgb]{0.56,0.35,0.01}{\textbf{\textit{#1}}}}
\newcommand{\ErrorTok}[1]{\textcolor[rgb]{0.64,0.00,0.00}{\textbf{#1}}}
\newcommand{\ExtensionTok}[1]{#1}
\newcommand{\FloatTok}[1]{\textcolor[rgb]{0.00,0.00,0.81}{#1}}
\newcommand{\FunctionTok}[1]{\textcolor[rgb]{0.13,0.29,0.53}{\textbf{#1}}}
\newcommand{\ImportTok}[1]{#1}
\newcommand{\InformationTok}[1]{\textcolor[rgb]{0.56,0.35,0.01}{\textbf{\textit{#1}}}}
\newcommand{\KeywordTok}[1]{\textcolor[rgb]{0.13,0.29,0.53}{\textbf{#1}}}
\newcommand{\NormalTok}[1]{#1}
\newcommand{\OperatorTok}[1]{\textcolor[rgb]{0.81,0.36,0.00}{\textbf{#1}}}
\newcommand{\OtherTok}[1]{\textcolor[rgb]{0.56,0.35,0.01}{#1}}
\newcommand{\PreprocessorTok}[1]{\textcolor[rgb]{0.56,0.35,0.01}{\textit{#1}}}
\newcommand{\RegionMarkerTok}[1]{#1}
\newcommand{\SpecialCharTok}[1]{\textcolor[rgb]{0.81,0.36,0.00}{\textbf{#1}}}
\newcommand{\SpecialStringTok}[1]{\textcolor[rgb]{0.31,0.60,0.02}{#1}}
\newcommand{\StringTok}[1]{\textcolor[rgb]{0.31,0.60,0.02}{#1}}
\newcommand{\VariableTok}[1]{\textcolor[rgb]{0.00,0.00,0.00}{#1}}
\newcommand{\VerbatimStringTok}[1]{\textcolor[rgb]{0.31,0.60,0.02}{#1}}
\newcommand{\WarningTok}[1]{\textcolor[rgb]{0.56,0.35,0.01}{\textbf{\textit{#1}}}}
\usepackage{graphicx}
\makeatletter
\newsavebox\pandoc@box
\newcommand*\pandocbounded[1]{% scales image to fit in text height/width
  \sbox\pandoc@box{#1}%
  \Gscale@div\@tempa{\textheight}{\dimexpr\ht\pandoc@box+\dp\pandoc@box\relax}%
  \Gscale@div\@tempb{\linewidth}{\wd\pandoc@box}%
  \ifdim\@tempb\p@<\@tempa\p@\let\@tempa\@tempb\fi% select the smaller of both
  \ifdim\@tempa\p@<\p@\scalebox{\@tempa}{\usebox\pandoc@box}%
  \else\usebox{\pandoc@box}%
  \fi%
}
% Set default figure placement to htbp
\def\fps@figure{htbp}
\makeatother
\setlength{\emergencystretch}{3em} % prevent overfull lines
\providecommand{\tightlist}{%
  \setlength{\itemsep}{0pt}\setlength{\parskip}{0pt}}
\usepackage{bookmark}
\IfFileExists{xurl.sty}{\usepackage{xurl}}{} % add URL line breaks if available
\urlstyle{same}
\hypersetup{
  pdftitle={Network Modeling: Assignment 2},
  hidelinks,
  pdfcreator={LaTeX via pandoc}}

\title{Network Modeling: Assignment 2}
\author{}
\date{\vspace{-2.5em}}

\begin{document}
\maketitle

\textbf{Group number}: 21\\
\textbf{Group members}: Eduard Mihai, Hanzhang Wang, Silvan
Affentranger, Linus Michel

\section{Libraries}\label{libraries}

\begin{Shaded}
\begin{Highlighting}[]
\CommentTok{\# Install packages from CRAN}
\FunctionTok{install.packages}\NormalTok{(}
\FunctionTok{c}\NormalTok{(}\StringTok{"network"}\NormalTok{, }\StringTok{"ergm"}\NormalTok{, }\StringTok{"manynet"}\NormalTok{, }\StringTok{"patchwork"}\NormalTok{, }\StringTok{"latticeExtra"}\NormalTok{, }\StringTok{"texreg"}\NormalTok{, }\StringTok{"broom"}\NormalTok{)}
\NormalTok{)}
\CommentTok{\# Install autograph from GitHub develop branch}
\CommentTok{\# install.packages("remotes") \# if not already installed}
\NormalTok{remotes}\SpecialCharTok{::}\FunctionTok{install\_github}\NormalTok{(}\StringTok{"stocnet/autograph"}\NormalTok{, }\AttributeTok{ref =} \StringTok{"develop"}\NormalTok{)}
\NormalTok{tinytex}\SpecialCharTok{::}\FunctionTok{install\_tinytex}\NormalTok{()}
\end{Highlighting}
\end{Shaded}

\begin{Shaded}
\begin{Highlighting}[]
\CommentTok{\# Loading libraries {-} stocnet suite}
\FunctionTok{library}\NormalTok{(manynet)}
\FunctionTok{library}\NormalTok{(autograph)}
\CommentTok{\# Loading libraries {-} sna suite}
\FunctionTok{library}\NormalTok{(network)}
\end{Highlighting}
\end{Shaded}

\begin{verbatim}
## 
## 'network' 1.20.0 (2026-02-06), part of the Statnet Project
## * 'news(package="network")' for changes since last version
## * 'citation("network")' for citation information
## * 'https://statnet.org' for help, support, and other information
\end{verbatim}

\begin{Shaded}
\begin{Highlighting}[]
\FunctionTok{library}\NormalTok{(ergm)}
\end{Highlighting}
\end{Shaded}

\begin{verbatim}
## 
## 'ergm' 4.12.0 (2026-02-17), part of the Statnet Project
## * 'news(package="ergm")' for changes since last version
## * 'citation("ergm")' for citation information
## * 'https://statnet.org' for help, support, and other information
\end{verbatim}

\begin{verbatim}
## 'ergm' 4 is a major update that introduces some backwards-incompatible
## changes. Please type 'news(package="ergm")' for a list of major
## changes.
\end{verbatim}

\begin{Shaded}
\begin{Highlighting}[]
\CommentTok{\# Loading libraries {-} plotting and presentation}
\FunctionTok{library}\NormalTok{(ggraph)}
\end{Highlighting}
\end{Shaded}

\begin{verbatim}
## Lade nötiges Paket: ggplot2
\end{verbatim}

\begin{Shaded}
\begin{Highlighting}[]
\FunctionTok{library}\NormalTok{(patchwork)}
\CommentTok{\# library(latticeExtra) \# improve the aesthetics of the MCMC diagnostics plots}
\FunctionTok{library}\NormalTok{(texreg)}
\end{Highlighting}
\end{Shaded}

\begin{verbatim}
## Version:  1.39.5
## Date:     2025-12-21
## Author:   Philip Leifeld (University of Manchester)
## 
## Consider submitting praise using the praise or praise_interactive functions.
## Please cite the JSS article in your publications -- see citation("texreg").
\end{verbatim}

\begin{Shaded}
\begin{Highlighting}[]
\FunctionTok{library}\NormalTok{(broom)}
\FunctionTok{library}\NormalTok{(sna)}
\end{Highlighting}
\end{Shaded}

\begin{verbatim}
## Lade nötiges Paket: statnet.common
\end{verbatim}

\begin{verbatim}
## 
## Attache Paket: 'statnet.common'
\end{verbatim}

\begin{verbatim}
## Die folgenden Objekte sind maskiert von 'package:base':
## 
##     attr, order, replace
\end{verbatim}

\begin{verbatim}
## sna: Tools for Social Network Analysis
## Version 2.8 created on 2024-09-07.
## copyright (c) 2005, Carter T. Butts, University of California-Irvine
##  For citation information, type citation("sna").
##  Type help(package="sna") to get started.
\end{verbatim}

\begin{Shaded}
\begin{Highlighting}[]
\FunctionTok{library}\NormalTok{(xtable)}
\FunctionTok{library}\NormalTok{(ggplot2)}
\FunctionTok{library}\NormalTok{(tidyr)}
\end{Highlighting}
\end{Shaded}

\begin{verbatim}
## 
## Attache Paket: 'tidyr'
\end{verbatim}

\begin{verbatim}
## Das folgende Objekt ist maskiert 'package:texreg':
## 
##     extract
\end{verbatim}

\section{Task 1: Geometrically Weighted Terms in
ERGMs}\label{task-1-geometrically-weighted-terms-in-ergms}

\subsection{1.1}\label{section}

\subsubsection{1.1.i}\label{i}

When a directed edge is added to a sender node with current out-degree
k, exactly two terms in the GWD statistic are affected: the count of
nodes with out-degree k decreases by one, and the count of nodes with
out-degree k+1 increases by one. Thus, if we take the difference
\(s^{\text{GWD}}(x', \alpha) - s^{\text{GWD}}(x, \alpha)\), all terms in
the sum except k and k+1 cancel out. Then we write out what is left
giving us
\(e^\alpha \left[ \left(1 - \left(1 - e^{-\alpha}\right)^{k+1}\right) - \left(1 - \left(1 - e^{-\alpha}\right)^{k}\right) \right]\)
and cancel the two ones:
\(e^\alpha \left[ \left(1 - e^{-\alpha}\right)^{k} - \left(1 - e^{-\alpha}\right)^{k+1} \right]\).
Next we factor out \(\left(1 - e^{-\alpha}\right) ^{k}\) leaving us with
\(e^\alpha \left(1 - e^{-\alpha}\right)^{k} \left[ 1 - \left(1 - e^{-\alpha}\right) \right]\).
Finally we realize \(1 - \left(1 - e^{-\alpha}\right) = e^{-\alpha}\)
and thus we get
\(\delta^{\text{GWD}}(k, \alpha) = \left(1 - e^{-\alpha}\right)^{k}\).

\begin{align*}
\delta^{\text{GWD}}(k, \alpha) &= s^{\text{GWD}}(x', \alpha) - s^{\text{GWD}}(x, \alpha) \\
&= e^\alpha \left[ \left(1 - \left(1 - e^{-\alpha}\right)^{k+1}\right) - \left(1 - \left(1 - e^{-\alpha}\right)^{k}\right) \right] \\
&= e^\alpha \left[ \left(1 - e^{-\alpha}\right)^{k} - \left(1 - e^{-\alpha}\right)^{k+1} \right] \\
&= e^\alpha \left(1 - e^{-\alpha}\right)^{k} \left[ 1 - \left(1 - e^{-\alpha}\right) \right] \\
&= e^\alpha \left(1 - e^{-\alpha}\right)^{k} \cdot e^{-\alpha} \\
&= \boxed{\left(1 - e^{-\alpha}\right)^{k}}
\end{align*}

\subsubsection{1.1.ii}\label{ii}

\begin{Shaded}
\begin{Highlighting}[]
\NormalTok{f1 }\OtherTok{\textless{}{-}} \FunctionTok{as.matrix}\NormalTok{(}\FunctionTok{read.csv}\NormalTok{(}\StringTok{"Glasgow/f1.csv"}\NormalTok{, }\AttributeTok{header =} \ConstantTok{FALSE}\NormalTok{))}
\NormalTok{f2 }\OtherTok{\textless{}{-}} \FunctionTok{as.matrix}\NormalTok{(}\FunctionTok{read.csv}\NormalTok{(}\StringTok{"Glasgow/f2.csv"}\NormalTok{, }\AttributeTok{header =} \ConstantTok{FALSE}\NormalTok{))}
\NormalTok{f3 }\OtherTok{\textless{}{-}} \FunctionTok{as.matrix}\NormalTok{(}\FunctionTok{read.csv}\NormalTok{(}\StringTok{"Glasgow/f3.csv"}\NormalTok{, }\AttributeTok{header =} \ConstantTok{FALSE}\NormalTok{))}

\NormalTok{demographic }\OtherTok{\textless{}{-}} \FunctionTok{read.csv}\NormalTok{(}\StringTok{"Glasgow/demographic.csv"}\NormalTok{)}
\NormalTok{logdistance }\OtherTok{\textless{}{-}} \FunctionTok{as.matrix}\NormalTok{(}\FunctionTok{read.csv}\NormalTok{(}\StringTok{"Glasgow/logdistance.csv"}\NormalTok{, }\AttributeTok{header =} \ConstantTok{FALSE}\NormalTok{))}
\NormalTok{alcohol }\OtherTok{\textless{}{-}} \FunctionTok{read.csv}\NormalTok{(}\StringTok{"Glasgow/alcohol.csv"}\NormalTok{)}

\CommentTok{\# out degree range from f1}
\NormalTok{max\_k }\OtherTok{\textless{}{-}} \FunctionTok{max}\NormalTok{(}\FunctionTok{rowSums}\NormalTok{(f1))}
\NormalTok{k\_vals }\OtherTok{\textless{}{-}} \DecValTok{0}\SpecialCharTok{:}\NormalTok{max\_k}

\CommentTok{\# change statistic function}
\NormalTok{change\_stat }\OtherTok{\textless{}{-}} \ControlFlowTok{function}\NormalTok{(k, a) \{}
  \ControlFlowTok{if}\NormalTok{ (a }\SpecialCharTok{==} \DecValTok{0}\NormalTok{) \{}
    \FunctionTok{return}\NormalTok{(}\FunctionTok{ifelse}\NormalTok{(k}\SpecialCharTok{==}\DecValTok{0}\NormalTok{, }\DecValTok{1}\NormalTok{, }\DecValTok{0}\NormalTok{))}
\NormalTok{  \}}
  \FunctionTok{return}\NormalTok{((}\DecValTok{1} \SpecialCharTok{{-}} \FunctionTok{exp}\NormalTok{(}\SpecialCharTok{{-}}\NormalTok{a))}\SpecialCharTok{\^{}}\NormalTok{k)}
\NormalTok{\}}

\CommentTok{\# alpha values to test}
\NormalTok{decay\_val\_neg }\OtherTok{\textless{}{-}} \FunctionTok{c}\NormalTok{(}\SpecialCharTok{{-}}\DecValTok{1}\NormalTok{, }\SpecialCharTok{{-}}\FloatTok{0.5}\NormalTok{)}
\NormalTok{decay\_val\_pos }\OtherTok{\textless{}{-}} \FunctionTok{c}\NormalTok{(}\DecValTok{0}\NormalTok{, }\FloatTok{0.1}\NormalTok{, }\FloatTok{0.2}\NormalTok{, }\FloatTok{0.3}\NormalTok{, }\FloatTok{0.5}\NormalTok{, }\FunctionTok{log}\NormalTok{(}\DecValTok{2}\NormalTok{), }\DecValTok{1}\NormalTok{, }\DecValTok{2}\NormalTok{, }\DecValTok{3}\NormalTok{)}

\CommentTok{\# generate plots for negative values of alpha}
\FunctionTok{par}\NormalTok{(}\AttributeTok{mfrow =} \FunctionTok{c}\NormalTok{(}\DecValTok{1}\NormalTok{,}\DecValTok{2}\NormalTok{))}
\ControlFlowTok{for}\NormalTok{ (a }\ControlFlowTok{in}\NormalTok{ decay\_val\_neg) \{}
\NormalTok{  y\_vals }\OtherTok{\textless{}{-}} \FunctionTok{sapply}\NormalTok{(k\_vals, change\_stat, }\AttributeTok{a =}\NormalTok{ a)}
  \FunctionTok{plot}\NormalTok{(k\_vals, y\_vals, }\AttributeTok{type =} \StringTok{"b"}\NormalTok{, }\AttributeTok{pch =} \DecValTok{19}\NormalTok{, }\AttributeTok{col =} \StringTok{"red"}\NormalTok{,}
       \AttributeTok{main =} \FunctionTok{paste}\NormalTok{(}\StringTok{"Alpha ="}\NormalTok{, a), }\AttributeTok{xlab =} \StringTok{"Out{-}degree (k)"}\NormalTok{, }\AttributeTok{ylab =} \StringTok{"Change Statistic"}\NormalTok{)}
  \FunctionTok{abline}\NormalTok{(}\AttributeTok{h =} \DecValTok{0}\NormalTok{, }\AttributeTok{lty =} \DecValTok{2}\NormalTok{)}
\NormalTok{\}}
\end{Highlighting}
\end{Shaded}

\pandocbounded{\includegraphics[keepaspectratio]{Submission2_files/Submission2_files/figure-latex/unnamed-chunk-4-1.pdf}}

\begin{Shaded}
\begin{Highlighting}[]
\CommentTok{\# generate plots for positive values of alpha}
\FunctionTok{par}\NormalTok{(}\AttributeTok{mfrow =} \FunctionTok{c}\NormalTok{(}\DecValTok{1}\NormalTok{, }\DecValTok{1}\NormalTok{))}
\NormalTok{colors }\OtherTok{\textless{}{-}} \FunctionTok{rainbow}\NormalTok{(}\FunctionTok{length}\NormalTok{(decay\_val\_pos))}
\FunctionTok{plot}\NormalTok{(}\ConstantTok{NULL}\NormalTok{, }\AttributeTok{xlim =} \FunctionTok{c}\NormalTok{(}\DecValTok{0}\NormalTok{, max\_k), }\AttributeTok{ylim =} \FunctionTok{c}\NormalTok{(}\DecValTok{0}\NormalTok{, }\DecValTok{1}\NormalTok{), }\AttributeTok{xlab =} \StringTok{"Out{-}degree (k)"}\NormalTok{,}
     \AttributeTok{ylab =} \StringTok{"Change Statistic"}\NormalTok{, }\AttributeTok{main =} \StringTok{"GWD Change Statistic (Non{-}negative Alpha)"}\NormalTok{)}

\ControlFlowTok{for}\NormalTok{ (i }\ControlFlowTok{in} \DecValTok{1}\SpecialCharTok{:}\FunctionTok{length}\NormalTok{(decay\_val\_pos)) \{}
\NormalTok{  y\_vals }\OtherTok{\textless{}{-}} \FunctionTok{sapply}\NormalTok{ (k\_vals, change\_stat, }\AttributeTok{a =}\NormalTok{ decay\_val\_pos[i])}
  \FunctionTok{lines}\NormalTok{(k\_vals, y\_vals, }\AttributeTok{type =} \StringTok{"b"}\NormalTok{, }\AttributeTok{pch =} \DecValTok{19}\NormalTok{, }\AttributeTok{col =}\NormalTok{ colors[i])}
\NormalTok{\}}

\FunctionTok{legend}\NormalTok{(}\StringTok{"topright"}\NormalTok{, }\AttributeTok{legend =} \FunctionTok{round}\NormalTok{(decay\_val\_pos, }\DecValTok{2}\NormalTok{), }\AttributeTok{col =}\NormalTok{ colors, }\AttributeTok{lty =} \DecValTok{1}\NormalTok{, }\AttributeTok{pch =} \DecValTok{19}\NormalTok{, }\AttributeTok{title =} \StringTok{"Alpha"}\NormalTok{, }\AttributeTok{cex =} \FloatTok{0.8}\NormalTok{)}
\end{Highlighting}
\end{Shaded}

\pandocbounded{\includegraphics[keepaspectratio]{Submission2_files/Submission2_files/figure-latex/unnamed-chunk-4-2.pdf}}

\subsubsection{1.1.iii}\label{iii}

The figures from the previous task show us that the decay parameter is
the primary influence on the diminishing marginal returns principle when
using gwodegree. If we look at the change statistics, we can get the
following takeaways, regarding how additional edges influence already
popular nodes. The first takeaway is that for all non-negative decay
values, the marginal contribution of an additional edge is strictly
non-increasing. As can be seen in our plot, the weight added to the
network statistic by a new edge is always maximized when the node where
the edge originates has an out-degree of zero. As k increases, this
contribution decreases, meaning that with each additional edge a node
has, the weights of that nodes connections decrease. With the
all-positive decay values, the value of alpha gives us the speed of this
decrease (as can be seen in the plot). With small values of alpha, the
change statistic drops to zero shortly after the first couple of edges,
while the marginal contribution of nodes with higher out-degrees stays
high with larger values of alpha. This allows us to choose how
influential we want high-degree hub-nodes to be on the overall network
statistic. The second takeaway is about the negative values for the
decay parameter alpha. For negative values it has a very unstable
behavior, as we can see in the way its graphs jump around. This is,
because the base of our change statistic becomes negative and thus
oscillates between positive and negative values, as the out-degree
increases. This means that if a second edge is added to a node, this
second tie would decrease the change statistic, while adding a third
such tie would increase it again. Applied to real world scenarios, this
makes little sense, which tells us that a positive alpha as decay
parameter value is needed to have meaningful results.

\subsubsection{1.1.iv}\label{iv}

To capture an activity effect where high out-degree nodes drive the
statistic, we need to choose the decay parameter alpha to be relatively
high. How high exactly depends on how high the out-degree nodes we want
to drive the statistic are. As we can see in the plot of positive values
for alpha, the choice of alpha influences the rate of diminishing
marginal returns for each additional outgoing tie. When alpha is small,
the change statistic drops to zero almost immediately after the first
tie. The model essentially only distinguishes between isolates and
non-isolates, which means that a small alpha would thus not effectively
capture the behavior of high-activity nodes because their additional
ties would contribute only a very small amount to the network statistic.
If we choose a larger alpha (like 2 or 3 for example), the curve of the
change statistic flattens, which causes the marginal contribution of an
extra tie to remain substantial even as a node's out-degree increases.
This means that for a student in our dataset for example, if he already
named four friends, naming the fifth and even a sixth friend still adds
significant weight to the network statistic. This prevents the model
from saturating early and maintains that additional ties from already
active keep having marginal impact, which gives us our desired outcome,
where high out-degree nodes drive the statistic.

\subsection{1.2}\label{section-1}

\subsubsection{1.2.i}\label{i-1}

\begin{Shaded}
\begin{Highlighting}[]
\NormalTok{f1\_matrix }\OtherTok{\textless{}{-}} \FunctionTok{as.matrix}\NormalTok{(}\FunctionTok{read.csv}\NormalTok{(}\StringTok{"Glasgow/f1.csv"}\NormalTok{, }\AttributeTok{header =} \ConstantTok{FALSE}\NormalTok{))}
\NormalTok{glasgow\_net }\OtherTok{\textless{}{-}} \FunctionTok{as.network}\NormalTok{(f1\_matrix, }\AttributeTok{directed =} \ConstantTok{TRUE}\NormalTok{)}

\NormalTok{max\_k }\OtherTok{\textless{}{-}} \FunctionTok{max}\NormalTok{(}\FunctionTok{rowSums}\NormalTok{(f1\_matrix))}

\NormalTok{m1 }\OtherTok{\textless{}{-}} \FunctionTok{ergm}\NormalTok{(glasgow\_net }\SpecialCharTok{\textasciitilde{}}\NormalTok{ edges)}
\end{Highlighting}
\end{Shaded}

\begin{verbatim}
## Starting maximum pseudolikelihood estimation (MPLE):
\end{verbatim}

\begin{verbatim}
## Obtaining the responsible dyads.
\end{verbatim}

\begin{verbatim}
## Evaluating the predictor and response matrix.
\end{verbatim}

\begin{verbatim}
## Maximizing the pseudolikelihood.
\end{verbatim}

\begin{verbatim}
## Finished MPLE.
\end{verbatim}

\begin{verbatim}
## Evaluating log-likelihood at the estimate.
\end{verbatim}

\subsubsection{1.2.ii}\label{ii-1}

\begin{Shaded}
\begin{Highlighting}[]
\NormalTok{m2\_neg }\OtherTok{\textless{}{-}} \FunctionTok{ergm}\NormalTok{(glasgow\_net }\SpecialCharTok{\textasciitilde{}}\NormalTok{ edges }\SpecialCharTok{+} \FunctionTok{offset}\NormalTok{(}\FunctionTok{gwodegree}\NormalTok{(}\FloatTok{0.3}\NormalTok{, }\AttributeTok{fixed =} \ConstantTok{TRUE}\NormalTok{)),}
               \AttributeTok{offset.coef =} \FunctionTok{c}\NormalTok{(}\SpecialCharTok{{-}}\DecValTok{2}\NormalTok{))}
\end{Highlighting}
\end{Shaded}

\begin{verbatim}
## Starting maximum pseudolikelihood estimation (MPLE):
\end{verbatim}

\begin{verbatim}
## Obtaining the responsible dyads.
\end{verbatim}

\begin{verbatim}
## Evaluating the predictor and response matrix.
\end{verbatim}

\begin{verbatim}
## Maximizing the pseudolikelihood.
\end{verbatim}

\begin{verbatim}
## Finished MPLE.
\end{verbatim}

\begin{verbatim}
## Starting Monte Carlo maximum likelihood estimation (MCMLE):
\end{verbatim}

\begin{verbatim}
## Iteration 1 of at most 60:
\end{verbatim}

\begin{verbatim}
## Warning: 'glpk' selected as the solver, but package 'Rglpk' is not available;
## falling back to 'lpSolveAPI'. This should be fine unless the sample size and/or
## the number of parameters is very big.
\end{verbatim}

\begin{verbatim}
## 1
\end{verbatim}

\begin{verbatim}
## Optimizing with step length 0.7817.
\end{verbatim}

\begin{verbatim}
## The log-likelihood improved by 2.7266.
\end{verbatim}

\begin{verbatim}
## Estimating equations are not within tolerance region.
\end{verbatim}

\begin{verbatim}
## Iteration 2 of at most 60:
\end{verbatim}

\begin{verbatim}
## 1 Optimizing with step length 1.0000.
## The log-likelihood improved by 0.2416.
## Estimating equations are not within tolerance region.
## Iteration 3 of at most 60:
## 1 Optimizing with step length 1.0000.
## The log-likelihood improved by 0.0009.
## Convergence test p-value: 0.0114. Not converged with 99% confidence; increasing sample size.
## Iteration 4 of at most 60:
## 1 Optimizing with step length 1.0000.
## The log-likelihood improved by 0.0057.
## Convergence test p-value: 0.0069. Converged with 99% confidence.
## Finished MCMLE.
## Evaluating log-likelihood at the estimate. Fitting the dyad-independent submodel...
## Bridging between the dyad-independent submodel and the full model...
## Setting up bridge sampling...
## Using 16 bridges: 1 2 3 4 5 6 7 8 9 10 11 12 13 14 15 16 .
## Bridging finished.
## 
## This model was fit using MCMC.  To examine model diagnostics and check
## for degeneracy, use the mcmc.diagnostics() function.
\end{verbatim}

\begin{Shaded}
\begin{Highlighting}[]
\NormalTok{m2\_pos }\OtherTok{\textless{}{-}} \FunctionTok{ergm}\NormalTok{(glasgow\_net }\SpecialCharTok{\textasciitilde{}}\NormalTok{ edges }\SpecialCharTok{+} \FunctionTok{offset}\NormalTok{(}\FunctionTok{gwodegree}\NormalTok{(}\FloatTok{0.3}\NormalTok{, }\AttributeTok{fixed =} \ConstantTok{TRUE}\NormalTok{)),}
               \AttributeTok{offset.coef =} \FunctionTok{c}\NormalTok{(}\DecValTok{2}\NormalTok{))}
\end{Highlighting}
\end{Shaded}

\begin{verbatim}
## Starting maximum pseudolikelihood estimation (MPLE):
## Obtaining the responsible dyads.
## Evaluating the predictor and response matrix.
## Maximizing the pseudolikelihood.
## Finished MPLE.
## Starting Monte Carlo maximum likelihood estimation (MCMLE):
## Iteration 1 of at most 60:
## 1 Optimizing with step length 0.6254.
## The log-likelihood improved by 2.3459.
## Estimating equations are not within tolerance region.
## Iteration 2 of at most 60:
## 1 Optimizing with step length 1.0000.
## The log-likelihood improved by 0.4839.
## Estimating equations are not within tolerance region.
## Iteration 3 of at most 60:
## 1 Optimizing with step length 1.0000.
## The log-likelihood improved by 0.0245.
## Convergence test p-value: 0.0016. Converged with 99% confidence.
## Finished MCMLE.
## Evaluating log-likelihood at the estimate. Fitting the dyad-independent submodel...
## Bridging between the dyad-independent submodel and the full model...
## Setting up bridge sampling...
## Using 16 bridges: 1 2 3 4 5 6 7 8 9 10 11 12 13 14 15 16 .
## Bridging finished.
## 
## This model was fit using MCMC.  To examine model diagnostics and check
## for degeneracy, use the mcmc.diagnostics() function.
\end{verbatim}

\subsubsection{1.2.iii}\label{iii-1}

\begin{Shaded}
\begin{Highlighting}[]
\NormalTok{sim\_out\_dist }\OtherTok{\textless{}{-}} \ControlFlowTok{function}\NormalTok{(formula, coefficients, }\AttributeTok{n\_sim =} \DecValTok{100}\NormalTok{, max\_val) \{}
\NormalTok{  sims }\OtherTok{\textless{}{-}}\NormalTok{ ergm}\SpecialCharTok{::}\FunctionTok{simulate\_formula}\NormalTok{(formula, }\AttributeTok{coef =}\NormalTok{ coefficients, }\AttributeTok{nsim =}\NormalTok{ n\_sim)}
  
\NormalTok{  dist\_list }\OtherTok{\textless{}{-}} \FunctionTok{lapply}\NormalTok{(sims, }\ControlFlowTok{function}\NormalTok{(x) \{}
\NormalTok{    deg\_counts }\OtherTok{\textless{}{-}} \FunctionTok{table}\NormalTok{(}\FunctionTok{factor}\NormalTok{(}\FunctionTok{rowSums}\NormalTok{(}\FunctionTok{as.matrix}\NormalTok{(x)), }\AttributeTok{levels =} \DecValTok{0}\SpecialCharTok{:}\NormalTok{max\_val))}
    \FunctionTok{return}\NormalTok{(}\FunctionTok{as.data.frame}\NormalTok{(}\FunctionTok{t}\NormalTok{(}\FunctionTok{as.matrix}\NormalTok{(deg\_counts))))}
\NormalTok{  \})}
  
  \FunctionTok{return}\NormalTok{(}\FunctionTok{do.call}\NormalTok{(rbind, dist\_list))}
\NormalTok{\}}

\FunctionTok{set.seed}\NormalTok{(}\DecValTok{123}\NormalTok{)}

\CommentTok{\# M1}
\NormalTok{sim\_m1\_df }\OtherTok{\textless{}{-}} \FunctionTok{sim\_out\_dist}\NormalTok{(glasgow\_net }\SpecialCharTok{\textasciitilde{}}\NormalTok{ edges, }
                          \FunctionTok{coef}\NormalTok{(m1), }\DecValTok{100}\NormalTok{, max\_k)}

\CommentTok{\# M2 negative gwd}
\NormalTok{sim\_m2\_neg\_df }\OtherTok{\textless{}{-}} \FunctionTok{sim\_out\_dist}\NormalTok{(glasgow\_net }\SpecialCharTok{\textasciitilde{}}\NormalTok{ edges }\SpecialCharTok{+} \FunctionTok{offset}\NormalTok{(}\FunctionTok{gwodegree}\NormalTok{(}\FloatTok{0.3}\NormalTok{, }\AttributeTok{fixed =} \ConstantTok{TRUE}\NormalTok{)), }
                              \FunctionTok{coef}\NormalTok{(m2\_neg), }\DecValTok{100}\NormalTok{, max\_k)}
\CommentTok{\# M2 positive gwd}
\NormalTok{sim\_m2\_pos\_df }\OtherTok{\textless{}{-}} \FunctionTok{sim\_out\_dist}\NormalTok{(glasgow\_net }\SpecialCharTok{\textasciitilde{}}\NormalTok{ edges }\SpecialCharTok{+} \FunctionTok{offset}\NormalTok{(}\FunctionTok{gwodegree}\NormalTok{(}\FloatTok{0.3}\NormalTok{, }\AttributeTok{fixed =} \ConstantTok{TRUE}\NormalTok{)), }
                              \FunctionTok{coef}\NormalTok{(m2\_pos), }\DecValTok{100}\NormalTok{, max\_k)}

\CommentTok{\# formatting for plot}
\NormalTok{process\_df }\OtherTok{\textless{}{-}} \ControlFlowTok{function}\NormalTok{(df, label) \{}
\NormalTok{  df }\SpecialCharTok{\%\textgreater{}\%}
    \FunctionTok{pivot\_longer}\NormalTok{(}\AttributeTok{cols =} \FunctionTok{everything}\NormalTok{(), }\AttributeTok{names\_to =} \StringTok{"degree"}\NormalTok{, }\AttributeTok{values\_to =} \StringTok{"node\_count"}\NormalTok{) }\SpecialCharTok{\%\textgreater{}\%}
    \FunctionTok{mutate}\NormalTok{(}\AttributeTok{model =}\NormalTok{ label, }\AttributeTok{degree =} \FunctionTok{as.numeric}\NormalTok{(degree))}
\NormalTok{\}}

\CommentTok{\# combine data}
\NormalTok{plot\_data }\OtherTok{\textless{}{-}} \FunctionTok{rbind}\NormalTok{(}
  \FunctionTok{process\_df}\NormalTok{(sim\_m1\_df, }\StringTok{"M1 (Edges Only)"}\NormalTok{),}
  \FunctionTok{process\_df}\NormalTok{(sim\_m2\_neg\_df, }\StringTok{"M2 (GWOD Negative)"}\NormalTok{),}
  \FunctionTok{process\_df}\NormalTok{(sim\_m2\_pos\_df, }\StringTok{"M2 (GWOD Positive)"}\NormalTok{)}
\NormalTok{)}

\CommentTok{\# generate plot}
\FunctionTok{ggplot}\NormalTok{(plot\_data, }\FunctionTok{aes}\NormalTok{(}\AttributeTok{x =} \FunctionTok{factor}\NormalTok{(degree), }\AttributeTok{y =}\NormalTok{ node\_count, }\AttributeTok{fill =}\NormalTok{ model)) }\SpecialCharTok{+}
  \FunctionTok{geom\_violin}\NormalTok{(}\AttributeTok{trim =} \ConstantTok{TRUE}\NormalTok{, }\AttributeTok{scale =} \StringTok{"width"}\NormalTok{, }\AttributeTok{alpha =} \FloatTok{0.85}\NormalTok{) }\SpecialCharTok{+}
  \FunctionTok{facet\_wrap}\NormalTok{(}\SpecialCharTok{\textasciitilde{}}\NormalTok{model) }\SpecialCharTok{+}
  \FunctionTok{labs}\NormalTok{(}
    \AttributeTok{title    =} \StringTok{"Out{-}Degree Distribution Comparison"}\NormalTok{,}
    \AttributeTok{subtitle =} \StringTok{"Decay = 0.3, 100 Simulations"}\NormalTok{,}
    \AttributeTok{x        =} \StringTok{"Out{-}degree (k)"}\NormalTok{,}
    \AttributeTok{y        =} \StringTok{"Number of Nodes"}
\NormalTok{  ) }\SpecialCharTok{+}
  \FunctionTok{theme\_minimal}\NormalTok{() }\SpecialCharTok{+}
  \FunctionTok{theme}\NormalTok{(}\AttributeTok{legend.position =} \StringTok{"none"}\NormalTok{)}
\end{Highlighting}
\end{Shaded}

\pandocbounded{\includegraphics[keepaspectratio]{Submission2_files/Submission2_files/figure-latex/unnamed-chunk-7-1.pdf}}

\subsubsection{1.2.iv}\label{iv-1}

As we can see in the plot, the gwodegree alters the structural
tendencies of a network compared to a random baseline network. By
comparing the three different configurations, we can observe how
different gwodegrees influence the degree variance among nodes. The
baseline is our M1 configuration, which represents a Bernoulli random
network where the probability of a tie is uniformly spread across all
node pairs. As we can see, the out-degree distribution follows a typical
bell-curve, that peaks at an out-degree of three. This indicates that
most nodes have a similar number of outgoing ties and the extremes,
nodes having zero or six ties, are infrequent. In M2 with a negative
gwodegree, we can observe that with a negative gwodegree the
distribution is flatter and more dispersed than in the random baseline.
This is, because the model penalizes higher out-degrees if it has a
negative gwodegree. Another obvious difference to the random baseline
is, that the extremes of having six and especially zero ties are much
more likely. If we take a look at M2 with a positive gwodegree, we can
see from the plot, that it increases the number of nodes with multiple
outgoing ties. In this configuration, the distribution becomes much
narrower and taller and is concentrated around out-degrees of two, three
and four. What can also easily be noticed is that the number of isolated
nodes with an out-degree of zero drops to almost zero.

\section{Task 2: Estimation and interpretation of an
ERGM}\label{task-2-estimation-and-interpretation-of-an-ergm}

\subsection{2.1}\label{section-2}

\begin{Shaded}
\begin{Highlighting}[]
\CommentTok{\# Reading the adjacency matrices using manynet::read\_matrix()}
\NormalTok{friendship }\OtherTok{\textless{}{-}} \FunctionTok{read\_matrix}\NormalTok{(}\StringTok{"Glasgow/f1.csv"}\NormalTok{, }\AttributeTok{header =} \ConstantTok{FALSE}\NormalTok{)}
\NormalTok{demographic }\OtherTok{\textless{}{-}} \FunctionTok{read.csv}\NormalTok{(}\StringTok{"Glasgow/demographic.csv"}\NormalTok{, }\AttributeTok{header =} \ConstantTok{TRUE}\NormalTok{)}

\NormalTok{net }\OtherTok{\textless{}{-}} \FunctionTok{network}\NormalTok{(}\FunctionTok{as.matrix}\NormalTok{(friendship), }\AttributeTok{directed =} \ConstantTok{TRUE}\NormalTok{)}
\NormalTok{net }\SpecialCharTok{\%v\%} \StringTok{"gender"} \OtherTok{\textless{}{-}} \FunctionTok{as.character}\NormalTok{(demographic[, }\StringTok{"gender"}\NormalTok{])}
\FunctionTok{set.seed}\NormalTok{(}\DecValTok{161}\NormalTok{)}
\CommentTok{\# this from the last assignment: permutations \textless{}{-} 5000}
\CommentTok{\#estimate ERGM with parameters of tie existence and if that tie is same gender}
\NormalTok{model }\OtherTok{\textless{}{-}} \FunctionTok{ergm}\NormalTok{(net }\SpecialCharTok{\textasciitilde{}}\NormalTok{ edges }\SpecialCharTok{+} \FunctionTok{nodematch}\NormalTok{(}\StringTok{"gender"}\NormalTok{), }\AttributeTok{control =} \FunctionTok{control.ergm}\NormalTok{(}\AttributeTok{seed =} \DecValTok{161}\NormalTok{))}
\end{Highlighting}
\end{Shaded}

\begin{verbatim}
## Starting maximum pseudolikelihood estimation (MPLE):
\end{verbatim}

\begin{verbatim}
## Obtaining the responsible dyads.
\end{verbatim}

\begin{verbatim}
## Evaluating the predictor and response matrix.
\end{verbatim}

\begin{verbatim}
## Maximizing the pseudolikelihood.
\end{verbatim}

\begin{verbatim}
## Finished MPLE.
\end{verbatim}

\begin{verbatim}
## Evaluating log-likelihood at the estimate.
\end{verbatim}

\begin{Shaded}
\begin{Highlighting}[]
\CommentTok{\#saveRDS(model, "model.rds")}
\FunctionTok{summary}\NormalTok{(model)}
\end{Highlighting}
\end{Shaded}

\begin{verbatim}
## Call:
## ergm(formula = net ~ edges + nodematch("gender"), control = control.ergm(seed = 161))
## 
## Maximum Likelihood Results:
## 
##                  Estimate Std. Error MCMC % z value Pr(>|z|)    
## edges             -5.0133     0.1365      0  -36.72   <1e-04 ***
## nodematch.gender   2.0098     0.1460      0   13.77   <1e-04 ***
## ---
## Signif. codes:  0 '***' 0.001 '**' 0.01 '*' 0.05 '.' 0.1 ' ' 1
## 
##      Null Deviance: 22890  on 16512  degrees of freedom
##  Residual Deviance:  3824  on 16510  degrees of freedom
##  
## AIC: 3828  BIC: 3843  (Smaller is better. MC Std. Err. = 0)
\end{verbatim}

\begin{Shaded}
\begin{Highlighting}[]
\CommentTok{\#get conditional probs {-}{-}{-}{-}{-}{-}}
\CommentTok{\#probability that a tie exists between two nodes}
\NormalTok{theta\_edges }\OtherTok{\textless{}{-}} \FunctionTok{coef}\NormalTok{(model)[}\StringTok{"edges"}\NormalTok{]}
\CommentTok{\#additional chance a tie exists if the nodes are same gender}
\NormalTok{theta\_match }\OtherTok{\textless{}{-}} \FunctionTok{coef}\NormalTok{(model)[}\StringTok{"nodematch.gender"}\NormalTok{]}

\CommentTok{\# Log{-}odds of a tie between two same{-}gender students}
\NormalTok{logodds\_same }\OtherTok{\textless{}{-}}\NormalTok{ theta\_edges }\SpecialCharTok{+}\NormalTok{ theta\_match}

\CommentTok{\# Converted to probability (exp / (1 + epx))}
\NormalTok{prob\_same }\OtherTok{\textless{}{-}} \FunctionTok{plogis}\NormalTok{(logodds\_same)}

\CommentTok{\# same for different gender tie (nodematch = 0)}
\NormalTok{logodds\_diff }\OtherTok{\textless{}{-}}\NormalTok{ theta\_edges}
\NormalTok{prob\_diff }\OtherTok{\textless{}{-}} \FunctionTok{plogis}\NormalTok{(logodds\_diff)}


\FunctionTok{cat}\NormalTok{(}\FunctionTok{sprintf}\NormalTok{(}\StringTok{"  P(tie | same gender)     :  \%.4f  (\%.1f \%\%)}\SpecialCharTok{\textbackslash{}n}\StringTok{"}\NormalTok{, prob\_same, prob\_same }\SpecialCharTok{*} \DecValTok{100}\NormalTok{))}
\end{Highlighting}
\end{Shaded}

\begin{verbatim}
##   P(tie | same gender)     :  0.0473  (4.7 %)
\end{verbatim}

\begin{Shaded}
\begin{Highlighting}[]
\FunctionTok{cat}\NormalTok{(}\FunctionTok{sprintf}\NormalTok{(}\StringTok{"  P(tie | diff gender)     :  \%.4f  (\%.1f \%\%)}\SpecialCharTok{\textbackslash{}n}\StringTok{"}\NormalTok{, prob\_diff, prob\_diff }\SpecialCharTok{*} \DecValTok{100}\NormalTok{))}
\end{Highlighting}
\end{Shaded}

\begin{verbatim}
##   P(tie | diff gender)     :  0.0066  (0.7 %)
\end{verbatim}

\begin{Shaded}
\begin{Highlighting}[]
\FunctionTok{cat}\NormalTok{(}\FunctionTok{sprintf}\NormalTok{(}\StringTok{"  Ratio Same{-}Diff gender   :  \%.4f  }\SpecialCharTok{\textbackslash{}n}\StringTok{"}\NormalTok{, prob\_same}\SpecialCharTok{/}\NormalTok{prob\_diff))}
\end{Highlighting}
\end{Shaded}

\begin{verbatim}
##   Ratio Same-Diff gender   :  7.1563
\end{verbatim}

Both results hold strong evidence that edges are both sparse
(\(\theta_{edges}=-5.0133\)) and homophillic
(\(\theta_{match}=2.0098\)), as \(p << 0.05\) in both cases. Although
there is only a probability of \(P_{same}= 4.7\%\) for a tie to exist
between two students of the same gender, it is still \(7.1563\) times
higher than the probability of a tie existing between two students of
opposite gender \(P_{diff}=0.7\%\). The huge difference in Null Deviance
to Residual Deviance (\(22890 \rightarrow 3824\)) also tells that the
tie structure is very dependent on the chosen parameters.

\subsection{2.2}\label{section-3}

To build the big ERGM model similar to 2.1:

\begin{Shaded}
\begin{Highlighting}[]
\NormalTok{alcohol }\OtherTok{\textless{}{-}} \FunctionTok{read.csv}\NormalTok{(}\StringTok{"Glasgow/alcohol.csv"}\NormalTok{, }\AttributeTok{header =} \ConstantTok{TRUE}\NormalTok{)}

\CommentTok{\# Attach just the first column}
\NormalTok{net }\SpecialCharTok{\%v\%} \StringTok{"alcohol"} \OtherTok{\textless{}{-}}\NormalTok{ alcohol[, }\DecValTok{1}\NormalTok{]}
\end{Highlighting}
\end{Shaded}

\begin{Shaded}
\begin{Highlighting}[]
\CommentTok{\#this takes a lil while:}
\NormalTok{big\_model }\OtherTok{\textless{}{-}} \FunctionTok{ergm}\NormalTok{(net }\SpecialCharTok{\textasciitilde{}}\NormalTok{ edges }\SpecialCharTok{+} \FunctionTok{nodematch}\NormalTok{(}\StringTok{"gender"}\NormalTok{) }\CommentTok{\#same baseline as 2.1}
               \SpecialCharTok{+}\NormalTok{ mutual }\CommentTok{\#(i) Reciprocity}
               \SpecialCharTok{+} \FunctionTok{gwesp}\NormalTok{(}\FloatTok{0.1}\NormalTok{, }\AttributeTok{fixed =} \ConstantTok{TRUE}\NormalTok{) }\CommentTok{\#(ii) Transitivity}
               
               \SpecialCharTok{+} \FunctionTok{gwodegree}\NormalTok{(}\FloatTok{0.1}\NormalTok{, }\AttributeTok{fixed =} \ConstantTok{TRUE}\NormalTok{)    }\CommentTok{\# controls for out{-}degree distribution}
               \SpecialCharTok{+} \FunctionTok{gwidegree}\NormalTok{(}\FloatTok{0.1}\NormalTok{, }\AttributeTok{fixed =} \ConstantTok{TRUE}\NormalTok{)    }\CommentTok{\# controls for in{-}degree distribution}
               
               \SpecialCharTok{+} \FunctionTok{nodeofactor}\NormalTok{(}\StringTok{"gender"}\NormalTok{, }\AttributeTok{base =} \DecValTok{2}\NormalTok{) }\CommentTok{\#(iii) Boys have higher node outdegree}
               \SpecialCharTok{+} \FunctionTok{nodeifactor}\NormalTok{(}\StringTok{"gender"}\NormalTok{, }\AttributeTok{base =} \DecValTok{2}\NormalTok{) }\CommentTok{\#(iv) Boys have higher node indegree}
               \SpecialCharTok{+} \FunctionTok{nodeocov}\NormalTok{(}\StringTok{"alcohol"}\NormalTok{) }\CommentTok{\#(v) high alcohol {-}\textgreater{} high node outdegree}
               \SpecialCharTok{+} \FunctionTok{nodeicov}\NormalTok{(}\StringTok{"alcohol"}\NormalTok{) }\CommentTok{\#(vi) high alcohol {-}\textgreater{} high node indegree}
               \SpecialCharTok{+} \FunctionTok{absdiff}\NormalTok{(}\StringTok{"alcohol"}\NormalTok{), }\CommentTok{\#(vii) alcohol homophily}
               \AttributeTok{control =} \FunctionTok{control.ergm}\NormalTok{(}\AttributeTok{seed =} \DecValTok{161}\NormalTok{, }\AttributeTok{MCMLE.maxit =} \DecValTok{40}\NormalTok{, }\AttributeTok{MCMC.samplesize =} \DecValTok{2048}\NormalTok{, }
                \AttributeTok{MCMC.burnin =} \DecValTok{10000}\NormalTok{, }\AttributeTok{MCMC.interval =} \DecValTok{500}\NormalTok{)}
\NormalTok{               )}
\CommentTok{\#summary(big\_model)}
\CommentTok{\#mcmc.diagnostics(big\_model)}
\CommentTok{\#saveRDS(big\_model, "big\_model.rds")}

\NormalTok{bigger\_model }\OtherTok{\textless{}{-}} \FunctionTok{ergm}\NormalTok{(net }\SpecialCharTok{\textasciitilde{}}\NormalTok{ edges }\SpecialCharTok{+} \FunctionTok{nodematch}\NormalTok{(}\StringTok{"gender"}\NormalTok{) }\CommentTok{\#same baseline as 2.1}
               \SpecialCharTok{+}\NormalTok{ mutual }\CommentTok{\#(i) Reciprocity}
               \SpecialCharTok{+} \FunctionTok{gwesp}\NormalTok{(}\FloatTok{0.1}\NormalTok{, }\AttributeTok{fixed =} \ConstantTok{TRUE}\NormalTok{) }\CommentTok{\#(ii) Transitivity}
               
               \SpecialCharTok{+} \FunctionTok{gwodegree}\NormalTok{(}\FloatTok{0.1}\NormalTok{, }\AttributeTok{fixed =} \ConstantTok{TRUE}\NormalTok{)    }\CommentTok{\# controls for out{-}degree distribution}
               \SpecialCharTok{+} \FunctionTok{gwidegree}\NormalTok{(}\FloatTok{0.1}\NormalTok{, }\AttributeTok{fixed =} \ConstantTok{TRUE}\NormalTok{)    }\CommentTok{\# controls for in{-}degree distribution}
               
               \SpecialCharTok{+} \FunctionTok{nodeofactor}\NormalTok{(}\StringTok{"gender"}\NormalTok{, }\AttributeTok{base =} \DecValTok{2}\NormalTok{) }\CommentTok{\#(iii) Boys have higher node outdegree}
               \SpecialCharTok{+} \FunctionTok{nodeifactor}\NormalTok{(}\StringTok{"gender"}\NormalTok{, }\AttributeTok{base =} \DecValTok{2}\NormalTok{) }\CommentTok{\#(iv) Boys have higher node indegree}
               \SpecialCharTok{+} \FunctionTok{nodeocov}\NormalTok{(}\StringTok{"alcohol"}\NormalTok{) }\CommentTok{\#(v) high alcohol {-}\textgreater{} high node outdegree}
               \SpecialCharTok{+} \FunctionTok{nodeicov}\NormalTok{(}\StringTok{"alcohol"}\NormalTok{) }\CommentTok{\#(vi) high alcohol {-}\textgreater{} high node indegree}
               \SpecialCharTok{+} \FunctionTok{absdiff}\NormalTok{(}\StringTok{"alcohol"}\NormalTok{), }\CommentTok{\#(vii) alcohol homophily}
               \AttributeTok{control =} \FunctionTok{control.ergm}\NormalTok{(}\AttributeTok{seed =} \DecValTok{161}\NormalTok{, }\AttributeTok{MCMC.samplesize =} \DecValTok{50000}\NormalTok{, }\AttributeTok{init =} \FunctionTok{coef}\NormalTok{(big\_model),}
                  \AttributeTok{MCMC.burnin =} \DecValTok{100000}\NormalTok{, }\AttributeTok{MCMC.interval =} \DecValTok{2000}\NormalTok{, }\AttributeTok{MCMLE.maxit =} \DecValTok{60}\NormalTok{)}
\NormalTok{               )}
\FunctionTok{saveRDS}\NormalTok{(bigger\_model, }\StringTok{"bigger\_model.rds"}\NormalTok{)}
\end{Highlighting}
\end{Shaded}

\begin{Shaded}
\begin{Highlighting}[]
\NormalTok{bigger\_model }\OtherTok{\textless{}{-}} \FunctionTok{readRDS}\NormalTok{(}\StringTok{"bigger\_model.rds"}\NormalTok{)}
\FunctionTok{summary}\NormalTok{(bigger\_model)}
\end{Highlighting}
\end{Shaded}

\begin{verbatim}
## Call:
## ergm(formula = net ~ edges + nodematch("gender") + mutual + gwesp(0.1, 
##     fixed = TRUE) + gwodegree(0.1, fixed = TRUE) + gwidegree(0.1, 
##     fixed = TRUE) + nodeofactor("gender", base = 2) + nodeifactor("gender", 
##     base = 2) + nodeocov("alcohol") + nodeicov("alcohol") + absdiff("alcohol"), 
##     control = control.ergm(seed = 161, MCMC.samplesize = 50000, 
##         init = coef(big_model), MCMC.burnin = 1e+05, MCMC.interval = 2000, 
##         MCMLE.maxit = 60))
## 
## Monte Carlo Maximum Likelihood Results:
## 
##                      Estimate Std. Error MCMC % z value Pr(>|z|)    
## edges                -5.95141    0.17513      0 -33.983  < 1e-04 ***
## nodematch.gender      0.85499    0.10496      0   8.146  < 1e-04 ***
## mutual                3.02297    0.19107      0  15.821  < 1e-04 ***
## gwesp.OTP.fixed.0.1   1.67667    0.10814      0  15.505  < 1e-04 ***
## gwodeg.fixed.0.1      1.84220    0.48991      0   3.760 0.000170 ***
## gwideg.fixed.0.1      1.60488    0.44954      0   3.570 0.000357 ***
## nodeofactor.gender.1 -0.10194    0.15031      0  -0.678 0.497628    
## nodeifactor.gender.1  0.03781    0.14949      0   0.253 0.800318    
## nodeocov.alcohol     -0.11891    0.05540      0  -2.147 0.031823 *  
## nodeicov.alcohol      0.15648    0.05472      0   2.860 0.004242 ** 
## absdiff.alcohol      -0.11072    0.03710      0  -2.985 0.002840 ** 
## ---
## Signif. codes:  0 '***' 0.001 '**' 0.01 '*' 0.05 '.' 0.1 ' ' 1
## 
##      Null Deviance: 22890  on 16512  degrees of freedom
##  Residual Deviance:  2845  on 16501  degrees of freedom
##  
## AIC: 2867  BIC: 2952  (Smaller is better. MC Std. Err. = 0.6422)
\end{verbatim}

\begin{Shaded}
\begin{Highlighting}[]
\FunctionTok{mcmc.diagnostics}\NormalTok{(bigger\_model)}
\end{Highlighting}
\end{Shaded}

\pandocbounded{\includegraphics[keepaspectratio]{Submission2_files/Submission2_files/figure-latex/unnamed-chunk-11-1.pdf}}
\pandocbounded{\includegraphics[keepaspectratio]{Submission2_files/Submission2_files/figure-latex/unnamed-chunk-11-2.pdf}}
\pandocbounded{\includegraphics[keepaspectratio]{Submission2_files/Submission2_files/figure-latex/unnamed-chunk-11-3.pdf}}
\pandocbounded{\includegraphics[keepaspectratio]{Submission2_files/Submission2_files/figure-latex/unnamed-chunk-11-4.pdf}}

\begin{verbatim}
## Sample statistics summary:
## 
## Iterations = 5002000:99994000
## Thinning interval = 24000 
## Number of chains = 1 
## Sample size per chain = 3959 
## 
## 1. Empirical mean and standard deviation for each variable,
##    plus standard error of the mean:
## 
##                          Mean      SD Naive SE Time-series SE
## edges                 2.27936  46.589  0.74045        1.63933
## nodematch.gender      2.28416  45.868  0.72898        1.56563
## mutual                0.91917  17.175  0.27297        0.60137
## gwesp.OTP.fixed.0.1   2.96813  55.530  0.88255        1.87951
## gwodeg.fixed.0.1      0.04214   2.406  0.03824        0.05035
## gwideg.fixed.0.1      0.02785   2.643  0.04201        0.05697
## nodeofactor.gender.1  0.30639  37.669  0.59868        1.26962
## nodeifactor.gender.1  0.37232  37.670  0.59869        1.27141
## nodeocov.alcohol      6.33721 116.779  1.85597        3.97827
## nodeicov.alcohol      6.54584 122.325  1.94411        4.12807
## absdiff.alcohol      -0.26623  50.475  0.80220        1.57395
## 
## 2. Quantiles for each variable:
## 
##                          2.5%     25%     50%    75%   97.5%
## edges                 -84.000 -31.000  0.0000 33.000  95.000
## nodematch.gender      -83.000 -30.000  0.0000 34.000  93.000
## mutual                -32.000 -11.000  0.0000 13.000  35.050
## gwesp.OTP.fixed.0.1  -100.557 -36.721  1.3579 40.279 114.761
## gwodeg.fixed.0.1       -5.247  -1.533  0.1990  1.781   4.279
## gwideg.fixed.0.1       -5.663  -1.702  0.2078  1.961   4.607
## nodeofactor.gender.1  -67.000 -27.000 -1.0000 25.000  78.000
## nodeifactor.gender.1  -67.000 -27.000 -1.0000 24.500  80.000
## nodeocov.alcohol     -211.050 -76.000  2.0000 82.000 244.000
## nodeicov.alcohol     -214.100 -78.000  3.0000 85.000 257.050
## absdiff.alcohol       -94.000 -35.000 -2.0000 32.500 102.000
## 
## 
## Are sample statistics significantly different from observed?
##                edges nodematch.gender    mutual gwesp.OTP.fixed.0.1
## diff.      2.2793635        2.2841627 0.9191715            2.968131
## test stat. 1.3904276        1.4589449 1.5284688            1.579201
## P-val.     0.1643991        0.1445803 0.1263962            0.114290
##            gwodeg.fixed.0.1 gwideg.fixed.0.1 nodeofactor.gender.1
## diff.            0.04214309       0.02784562            0.3063905
## test stat.       0.83697626       0.48881339            0.2413251
## P-val.           0.40260591       0.62497381            0.8093032
##            nodeifactor.gender.1 nodeocov.alcohol nodeicov.alcohol
## diff.                 0.3723162         6.337206         6.545845
## test stat.            0.2928362         1.592956         1.585693
## P-val.                0.7696473         0.111170         0.112809
##            absdiff.alcohol       (Omni)
## diff.           -0.2662288           NA
## test stat.      -0.1691465 25.556528815
## P-val.           0.8656814  0.008052915
## 
## Sample statistics cross-correlations:
##                          edges nodematch.gender    mutual gwesp.OTP.fixed.0.1
## edges                1.0000000        0.9757230 0.9317709           0.9785293
## nodematch.gender     0.9757230        1.0000000 0.9258228           0.9679443
## mutual               0.9317709        0.9258228 1.0000000           0.9492862
## gwesp.OTP.fixed.0.1  0.9785293        0.9679443 0.9492862           1.0000000
## gwodeg.fixed.0.1     0.4708714        0.4498849 0.4483660           0.4266192
## gwideg.fixed.0.1     0.4648140        0.4462667 0.4432268           0.4142682
## nodeofactor.gender.1 0.8079867        0.7910742 0.7491889           0.7918747
## nodeifactor.gender.1 0.8032780        0.7862700 0.7437564           0.7878244
## nodeocov.alcohol     0.9591899        0.9360376 0.9016265           0.9416844
## nodeicov.alcohol     0.9589828        0.9346047 0.8913238           0.9387741
## absdiff.alcohol      0.8244945        0.8000643 0.7592274           0.7988675
##                      gwodeg.fixed.0.1 gwideg.fixed.0.1 nodeofactor.gender.1
## edges                       0.4708714        0.4648140            0.8079867
## nodematch.gender            0.4498849        0.4462667            0.7910742
## mutual                      0.4483660        0.4432268            0.7491889
## gwesp.OTP.fixed.0.1         0.4266192        0.4142682            0.7918747
## gwodeg.fixed.0.1            1.0000000        0.3386004            0.3433523
## gwideg.fixed.0.1            0.3386004        1.0000000            0.3411299
## nodeofactor.gender.1        0.3433523        0.3411299            1.0000000
## nodeifactor.gender.1        0.3418181        0.3348452            0.9838166
## nodeocov.alcohol            0.4740957        0.4333060            0.6935012
## nodeicov.alcohol            0.4713878        0.4260358            0.6956318
## absdiff.alcohol             0.4231015        0.4100971            0.5934806
##                      nodeifactor.gender.1 nodeocov.alcohol nodeicov.alcohol
## edges                           0.8032780        0.9591899        0.9589828
## nodematch.gender                0.7862700        0.9360376        0.9346047
## mutual                          0.7437564        0.9016265        0.8913238
## gwesp.OTP.fixed.0.1             0.7878244        0.9416844        0.9387741
## gwodeg.fixed.0.1                0.3418181        0.4740957        0.4713878
## gwideg.fixed.0.1                0.3348452        0.4333060        0.4260358
## nodeofactor.gender.1            0.9838166        0.6935012        0.6956318
## nodeifactor.gender.1            1.0000000        0.6904251        0.6871942
## nodeocov.alcohol                0.6904251        1.0000000        0.9858228
## nodeicov.alcohol                0.6871942        0.9858228        1.0000000
## absdiff.alcohol                 0.5893233        0.8225123        0.8360954
##                      absdiff.alcohol
## edges                      0.8244945
## nodematch.gender           0.8000643
## mutual                     0.7592274
## gwesp.OTP.fixed.0.1        0.7988675
## gwodeg.fixed.0.1           0.4231015
## gwideg.fixed.0.1           0.4100971
## nodeofactor.gender.1       0.5934806
## nodeifactor.gender.1       0.5893233
## nodeocov.alcohol           0.8225123
## nodeicov.alcohol           0.8360954
## absdiff.alcohol            1.0000000
## 
## Sample statistics auto-correlation:
## Chain 1 
##                edges nodematch.gender    mutual gwesp.OTP.fixed.0.1
## Lag 0      1.0000000        1.0000000 1.0000000           1.0000000
## Lag 24000  0.6202566        0.6435884 0.6289399           0.6386142
## Lag 48000  0.4116515        0.4228933 0.4054584           0.4197581
## Lag 72000  0.2867670        0.2873365 0.2813267           0.2894669
## Lag 96000  0.2025114        0.2028105 0.1939781           0.2016381
## Lag 120000 0.1355932        0.1311716 0.1288809           0.1262960
##            gwodeg.fixed.0.1 gwideg.fixed.0.1 nodeofactor.gender.1
## Lag 0            1.00000000       1.00000000            1.0000000
## Lag 24000        0.13820566       0.13443423            0.6361161
## Lag 48000        0.08097046       0.08423023            0.4178025
## Lag 72000        0.04893458       0.08089974            0.2817434
## Lag 96000        0.02195496       0.06069668            0.1961316
## Lag 120000       0.04307426       0.02231194            0.1336774
##            nodeifactor.gender.1 nodeocov.alcohol nodeicov.alcohol
## Lag 0                 1.0000000        1.0000000        1.0000000
## Lag 24000             0.6369396        0.6133477        0.5986978
## Lag 48000             0.4178373        0.4060623        0.3979328
## Lag 72000             0.2824987        0.2792334        0.2747402
## Lag 96000             0.1989718        0.1969080        0.1917472
## Lag 120000            0.1380993        0.1318763        0.1286975
##            absdiff.alcohol
## Lag 0            1.0000000
## Lag 24000        0.4868468
## Lag 48000        0.2910116
## Lag 72000        0.1863064
## Lag 96000        0.1473820
## Lag 120000       0.1101644
## 
## Sample statistics burn-in diagnostic (Geweke):
## Chain 1 
## 
## Fraction in 1st window = 0.1
## Fraction in 2nd window = 0.5 
## 
##                edges     nodematch.gender               mutual 
##          -0.15138298          -0.14588321          -0.06055932 
##  gwesp.OTP.fixed.0.1     gwodeg.fixed.0.1     gwideg.fixed.0.1 
##          -0.08622441           0.41516841           0.76217049 
## nodeofactor.gender.1 nodeifactor.gender.1     nodeocov.alcohol 
##           0.65095168           0.55865531          -0.10720981 
##     nodeicov.alcohol      absdiff.alcohol 
##          -0.25959011          -0.94812182 
## 
## Individual P-values (lower = worse):
##                edges     nodematch.gender               mutual 
##            0.8796736            0.8840136            0.9517102 
##  gwesp.OTP.fixed.0.1     gwodeg.fixed.0.1     gwideg.fixed.0.1 
##            0.9312880            0.6780186            0.4459583 
## nodeofactor.gender.1 nodeifactor.gender.1     nodeocov.alcohol 
##            0.5150777            0.5763970            0.9146225 
##     nodeicov.alcohol      absdiff.alcohol 
##            0.7951800            0.3430674 
## Joint P-value (lower = worse):  0.2120706 
## 
## Note: To save space, only one in every 12 iterations of the MCMC sample
##   used for estimation was stored for diagnostics. Sample size per chain
##   was originally around 47508 with thinning interval 2000.
## 
## Note: MCMC diagnostics shown here are from the last round of
##   simulation, prior to computation of final parameter estimates.
##   Because the final estimates are refinements of those used for this
##   simulation run, these diagnostics may understate model performance.
##   To directly assess the performance of the final model on in-model
##   statistics, please use the GOF command: gof(ergmFitObject,
##   GOF=~model).
\end{verbatim}

\subsection{2.3}\label{section-4}

The ERGM model here itself has the 7 parameters to check for the wanted
features (\(i\) - \(vii\)). To achieve convergence we had to add in 2
more control parameters: \(gwodegree\) and \(gwidegree\). We see in the
summary, that both \(gwodegree\) and \(gwidegree\) are significant, so
they were important to explain the network and help with the convergence
of the algorithm. Using a high decay value like \(\alpha = 0.3\) or
\(\alpha = 0.6\) gave problems while converging so we went for one on
the lower bound of \(\alpha = 0.1\), since the MCMC sampler tended to
generate networks with way more edges than observed. The low \(\alpha\)
decay value also tells that in the network, transitivity is only
important over short distances, and longer chains do not contribute
much. Using small parameters in the control parameter, the
\(mcmc.diagnostics()\) gave some bad results: The Mean was far away from
0 for many terms, the P-value was below or very close to \(0.05\) and
the omnibus value being at \(3.4*10^{-6}\) indicates the model has not
fully converged yet. All of this indicated that a longer MCMC run with
larger parameters was needed as it is now the case. So from this first
model we now expand on it further by changing the parameters in control
to larger samplesize, burnin and interval. Further all this is done on
the initial model as it just has not fully converged yet, and starting
from these large settings gave a failure of convergence, as these are
bad starting parameters.

This second ERGM fit may have only one iteration until it finishes, but
it improves most statistics shown in \(mcmc.diagnostics()\). All terms
now have a P-value above \(0.05\), but most are still not close to 1.
All individual Geweke statistics are within acceptable bounds and the
Geweke joint P-value improved from \(0.089\) to \(0.212\), confirming
chain stationarity. The omnibus P-value is significantly bigger at
\(0.008>>3.4*10^{-6}\) but still below \(0.05\), which is fine since it
is sensitive to large sample sizes and not indicative of systematic
misfit in any specific statistic.

\subsection{2.4}\label{section-5}

\begin{Shaded}
\begin{Highlighting}[]
\NormalTok{gof\_model }\OtherTok{\textless{}{-}} \FunctionTok{gof}\NormalTok{(bigger\_model, }\AttributeTok{GOF =} \SpecialCharTok{\textasciitilde{}}\NormalTok{ idegree }\SpecialCharTok{+}\NormalTok{ odegree }\SpecialCharTok{+}\NormalTok{ espartners }\SpecialCharTok{+}\NormalTok{ distance) }\CommentTok{\#+ dspartners}
\FunctionTok{saveRDS}\NormalTok{(gof\_model, }\StringTok{"gof\_model.rds"}\NormalTok{)}
\end{Highlighting}
\end{Shaded}

\begin{Shaded}
\begin{Highlighting}[]
\NormalTok{gof\_model }\OtherTok{\textless{}{-}} \FunctionTok{readRDS}\NormalTok{(}\StringTok{"gof\_model.rds"}\NormalTok{)}
\CommentTok{\# Plot auxiliary statistics}
\NormalTok{gof\_model}
\end{Highlighting}
\end{Shaded}

\begin{verbatim}
## 
## Goodness-of-fit for model statistics 
## 
##                            obs      min      mean       max MC p-value
## edges                 448.0000 345.0000  447.1000  554.0000       0.94
## nodematch.gender      394.0000 281.0000  393.6400  504.0000       0.88
## mutual                124.0000  84.0000  125.8500  163.0000       0.92
## gwesp.OTP.fixed.0.1   325.1830 200.0814  325.3779  460.9449       0.98
## gwodeg.fixed.0.1      134.9205 125.6260  134.8538  140.3149       0.96
## gwideg.fixed.0.1      133.7391 125.7896  133.8034  140.1647       1.00
## nodeofactor.gender.1  266.0000 181.0000  265.2500  360.0000       1.00
## nodeifactor.gender.1  266.0000 185.0000  265.7900  358.0000       0.98
## nodeocov.alcohol     1084.0000 787.0000 1087.0400 1377.0000       1.00
## nodeicov.alcohol     1136.0000 816.0000 1139.7600 1461.0000       0.92
## absdiff.alcohol       446.0000 334.0000  444.7400  567.0000       0.94
## 
## Goodness-of-fit for minimum geodesic distance 
## 
##      obs  min    mean  max MC p-value
## 1    448  345  447.10  554       0.94
## 2    777  834 1312.68 3677       0.00
## 3   1162 1424 2562.32 4118       0.00
## 4   1580 1737 3365.75 4564       0.00
## 5   1802 1818 3013.10 3770       0.00
## 6   1739  915 1927.51 2721       0.54
## 7   1625  164  986.96 1656       0.04
## 8   1274    8  455.64 1326       0.02
## 9    872    0  203.26  933       0.02
## 10   573    0   89.62  557       0.00
## 11   287    0   38.04  270       0.00
## 12   103    0   15.31  132       0.10
## 13    21    0    6.09   73       0.20
## 14     5    0    2.36   35       0.24
## 15     0    0    0.76   21       1.00
## 16     0    0    0.24    9       1.00
## 17     0    0    0.04    2       1.00
## Inf 4244  632 2085.22 4353       0.02
## 
## Goodness-of-fit for out-degree 
## 
##    obs min  mean max MC p-value
## 0    6   0  4.48  12       0.56
## 1    7  12 23.48  41       0.00
## 2   25  14 26.88  39       0.80
## 3   28  10 23.05  34       0.36
## 4   25  10 16.90  26       0.08
## 5   21   4 12.01  22       0.04
## 6   17   1  7.94  15       0.00
## 7    0   1  5.16  14       0.00
## 8    0   0  3.18   8       0.06
## 9    0   0  2.27   8       0.22
## 10   0   0  1.28   5       0.66
## 11   0   0  0.88   4       0.78
## 12   0   0  0.53   3       1.00
## 13   0   0  0.34   4       1.00
## 14   0   0  0.16   2       1.00
## 15   0   0  0.16   1       1.00
## 16   0   0  0.14   2       1.00
## 17   0   0  0.07   2       1.00
## 18   0   0  0.05   1       1.00
## 19   0   0  0.01   1       1.00
## 22   0   0  0.01   1       1.00
## 23   0   0  0.01   1       1.00
## 56   0   0  0.01   1       1.00
## 
## Goodness-of-fit for in-degree 
## 
##    obs min  mean max MC p-value
## 0    6   0  5.45  12       0.90
## 1   18  11 23.36  34       0.38
## 2   28  16 26.04  38       0.72
## 3   23  12 22.46  33       1.00
## 4   18   9 17.34  25       1.00
## 5   12   4 12.02  21       1.00
## 6    9   2  7.67  14       0.78
## 7    8   1  5.48  12       0.42
## 8    2   0  3.18   9       0.76
## 9    4   0  2.14   6       0.34
## 10   0   0  1.46   5       0.46
## 11   1   0  0.85   4       1.00
## 12   0   0  0.70   3       1.00
## 13   0   0  0.31   2       1.00
## 14   0   0  0.13   1       1.00
## 15   0   0  0.12   2       1.00
## 16   0   0  0.08   1       1.00
## 17   0   0  0.08   1       1.00
## 18   0   0  0.04   1       1.00
## 19   0   0  0.01   1       1.00
## 20   0   0  0.02   1       1.00
## 22   0   0  0.01   1       1.00
## 23   0   0  0.02   1       1.00
## 24   0   0  0.01   1       1.00
## 26   0   0  0.01   1       1.00
## 61   0   0  0.01   1       1.00
## 
## Goodness-of-fit for edgewise shared partner 
## 
##       obs min   mean max MC p-value
## .OTP0 137  92 130.07 162       0.50
## .OTP1 167 150 231.11 304       0.06
## .OTP2  92  28  67.39 110       0.28
## .OTP3  42   3  14.30  36       0.00
## .OTP4   9   0   3.26  17       0.18
## .OTP5   1   0   0.70  11       0.70
## .OTP6   0   0   0.22   5       1.00
## .OTP7   0   0   0.04   3       1.00
## .OTP9   0   0   0.01   1       1.00
\end{verbatim}

\begin{Shaded}
\begin{Highlighting}[]
\NormalTok{ergm}\SpecialCharTok{:::}\FunctionTok{plot.gof}\NormalTok{(gof\_model)}
\end{Highlighting}
\end{Shaded}

\pandocbounded{\includegraphics[keepaspectratio]{Submission2_files/Submission2_files/figure-latex/unnamed-chunk-13-1.pdf}}
\pandocbounded{\includegraphics[keepaspectratio]{Submission2_files/Submission2_files/figure-latex/unnamed-chunk-13-2.pdf}}
\pandocbounded{\includegraphics[keepaspectratio]{Submission2_files/Submission2_files/figure-latex/unnamed-chunk-13-3.pdf}}
\pandocbounded{\includegraphics[keepaspectratio]{Submission2_files/Submission2_files/figure-latex/unnamed-chunk-13-4.pdf}}
\pandocbounded{\includegraphics[keepaspectratio]{Submission2_files/Submission2_files/figure-latex/unnamed-chunk-13-5.pdf}}

Both the plot of \(model\) \(statistic\) and \(idegree\) indicate a very
good fit on these terms, as there are no outliers and their values are
very close to the centre of the boxplots. This confirms that the model
is internally consistent and that it reproduces the in-degree
distribution well, it accurately captures how many nominations students
receive. The \(odegree\) plot is somewhat predictably wrong, it
overestimates the count of nodes that have exactly \(1\) outgoing edge
and then has too few with those that have \(3\)-\(6\) edges. It wrongly
assumes that there are some nodes with more than \(6\) outgoing edges,
which is not possible per the data collection, as that only allowed up
to \(6\) friendship nominations per person, which is information that
the model does not know about. The \(edge\)-\(wise\) \(shared\)
\(partners\) plot shows that the model has problems with predicting
triads, which suggests that the chosen decay value \(\alpha = 0.1\)
might not have been the best choice. The bad fitting of the \(minimum\)
\(goedesic\) \(distance\) plot shows that the estimated model is more
fragmented and less connected than the real observed data would suggest.
Overall, the model shows clear parts that can still be improved, for one
there is the \(\alpha\) decay value that could be tweaked until there
are better results. There could also be more parameters that the model
does not count for yet, like the fact that a node has a maximum
out-degree of 6. It seems to be better with replicating local features,
while its weaknesses lie with the global features.

\subsection{2.5}\label{section-6}

\begin{Shaded}
\begin{Highlighting}[]
\FunctionTok{summary}\NormalTok{(bigger\_model)}
\end{Highlighting}
\end{Shaded}

\begin{verbatim}
## Call:
## ergm(formula = net ~ edges + nodematch("gender") + mutual + gwesp(0.1, 
##     fixed = TRUE) + gwodegree(0.1, fixed = TRUE) + gwidegree(0.1, 
##     fixed = TRUE) + nodeofactor("gender", base = 2) + nodeifactor("gender", 
##     base = 2) + nodeocov("alcohol") + nodeicov("alcohol") + absdiff("alcohol"), 
##     control = control.ergm(seed = 161, MCMC.samplesize = 50000, 
##         init = coef(big_model), MCMC.burnin = 1e+05, MCMC.interval = 2000, 
##         MCMLE.maxit = 60))
## 
## Monte Carlo Maximum Likelihood Results:
## 
##                      Estimate Std. Error MCMC % z value Pr(>|z|)    
## edges                -5.95141    0.17513      0 -33.983  < 1e-04 ***
## nodematch.gender      0.85499    0.10496      0   8.146  < 1e-04 ***
## mutual                3.02297    0.19107      0  15.821  < 1e-04 ***
## gwesp.OTP.fixed.0.1   1.67667    0.10814      0  15.505  < 1e-04 ***
## gwodeg.fixed.0.1      1.84220    0.48991      0   3.760 0.000170 ***
## gwideg.fixed.0.1      1.60488    0.44954      0   3.570 0.000357 ***
## nodeofactor.gender.1 -0.10194    0.15031      0  -0.678 0.497628    
## nodeifactor.gender.1  0.03781    0.14949      0   0.253 0.800318    
## nodeocov.alcohol     -0.11891    0.05540      0  -2.147 0.031823 *  
## nodeicov.alcohol      0.15648    0.05472      0   2.860 0.004242 ** 
## absdiff.alcohol      -0.11072    0.03710      0  -2.985 0.002840 ** 
## ---
## Signif. codes:  0 '***' 0.001 '**' 0.01 '*' 0.05 '.' 0.1 ' ' 1
## 
##      Null Deviance: 22890  on 16512  degrees of freedom
##  Residual Deviance:  2845  on 16501  degrees of freedom
##  
## AIC: 2867  BIC: 2952  (Smaller is better. MC Std. Err. = 0.6422)
\end{verbatim}

Terms not used for any hypothesis: The \(edges\) coefficient of
(\(-5.95, p<0.001\)) confirms that the friendship network is very sparse
and a the probability for a tie existing between any two students is
very low. Very logical for a huge network, where any student can only
nominate a maximum of \(6\) friends. Both \(gwodegree\)
(\(+1.84, p<0.001\)) and \(gwidegree\) (\(+1.60, p<0.001\)) are
significant, confirming there is a clear difference between how many
nominations students both send and receive. Some students are far more
social or popular than others. The gender homophily from 2.1 still stays
relevant here with \(nodematch.gender\) (\(+0.85, p<0.001\)) and
indicates that it plays an important role in how ties are formed. It is
to be noted though, that this term dropped from the model in 2.1 with a
value of \(+2.01\), so some of the networks model can be better
explained by other terms introduced in this new bigger model.

Hypothesis (\(i\)) Reciprocity: \(mutual\) (\(+3.02, p < 0.001\)). The
hypothesis is strongly supported by these values. If student j has
nominated student i, the odds of i nominating j back are multiplied by
\(e^{3.02} ≈ 20.5\). This is the strongest effect that is found in this
model, indicating it is a crucial deciding factor in how nominations are
built up on.

Hypothesis (\(ii\)) Transitivity: \(gwesp\) (\(+1.68, p < 0.001\)). The
hypothesis is strongly supported by these values. When two students
share a common friend, a tie between them becomes significantly more
likely. This supports a typical triadic closure structure in the model,
which occurs in most social networks and also leads to the formation of
clusters.

Hypothesis (\(iii\)) Boys have higher social activity:
\(nodeofactor.gender\) (\(−0.10, p = 0.498\)). This hypothesis is not
supported by these values. The coefficient for boys sending more
nominations than girls is small and negative, which would suggest that
the opposite is slightly true. This however is not the case either,
since the P-value is so high and as such, not making it statistically
significant.

Hypothesis (\(iv\)) Boys have higher popularity: \(nodeifactor.gender\)
(\(+0.04, p = 0.800\)). This hypothesis is not supported by these
values. The coefficient is near zero and the high P-value means it is
completely insignificant. Boys do not receive significantly more
nominations than girls in this network, which, along with Hypothesis
(\(iii\)), means that gender has no direct influence on how popular or
socially active a person is.

Hypothesis (\(v\)) Alcohol consumption increases social activity:
\(nodeocov.alcohol\) (\(−0.12, p = 0.032\)). This hypothesis is not
supported by these values, the somewhat low P-value along a negative
coefficient say the opposite is supported. Higher alcohol consumption is
associated with fewer outgoing nominations. A one unit increase in
alcohol score decreases the odds-ratio of sending a tie by
\(e^{0.12} ≈ 1.127\), suggesting that heavier drinkers are actually less
likely to actively nominate friends, contrary to the hypothesis. With
four units difference from \(1\) to \(5\), this means an approximate
odd-ratio decrease of \(e^{0.12 * 4} ≈ 1.616\) for the heaviest drinkers
in comparison to those with a value of \(1\).

Hypothesis (\(vi\)) Alcohol consumption increases popularity:
\(nodeicov.alcohol\) (\(+0.16, p = 0.004\)). The hypothesis is supported
by these values. Students with higher alcohol consumption receive
significantly more friendship nominations. A one unit increase in
alcohol score multiplies the odds of receiving a tie by
\(e^{0.16} ≈ 1.174\), suggesting heavier drinkers are perceived as more
popular or socially attractive by their peers. With four units
difference from \(1\) to \(5\), this means an approximate odd increase
of \(e^{0.16 * 4} ≈ 1.896\) for the heaviest drinkers in comparison to
those with a value of \(1\).

Hypothesis (\(vii\)) Alcohol homophily: \(absdiff.alcohol\)
(\(−0.11, p = 0.003\)). The hypothesis is supported by these values. The
negative coefficient on absolute difference means ties are more likely
between students with similar alcohol consumption levels. The smaller
the difference in alcohol scores between two students, the more likely
they are to form a friendship, indicating clear homophillic tendencies
regarding alcohol consumption.

\subsection{2.6}\label{section-7}

Since the data collection method only allowed the students to nominate 6
friends, a constraint that enforces this could be added to the ERGM
estimation model to get this point correct, which would make the model
itself more related to the model and fit some terms better. This would
however raise the question if this is really wanted, since it is very
possible that in real-life, students have more than 6 friendship
nominations. So this should rather a problem to think about at the data
collection part of data modelling. From the results we found in 2.4, we
should tweak the \(\alpha\) decay value to a higher value like 0.3 or
maybe even higher than that, since transitivity between students
friendships seems to hold over more nodes than the current value
estimates for and it does not correctly simulate the triadic patterns
that can be found in this network. There is a very arguable point to be
made that ties are also dependent on age-homophily, something that the
current model does not account for at all, since the edges do not even
have the age value added to them. There is really no reason not to add
the age parameter to the model, since we already have the data
available. Adding these terms:

\begin{Shaded}
\begin{Highlighting}[]
\CommentTok{\# + nodeocov("age") + nodeicov("age") + absdiff("age")}
\end{Highlighting}
\end{Shaded}

could help the model. Statistical outliers, like people who have to
repeat a class, may be on one end of the spectrum and more likely to
have drastically different data values than other students.

\subsection{2.7}\label{section-8}

We suggest two hypotheses that center around the uncollected data of the
students grades. The first hypothesis is centered around Academic
homophily. This would require the following addition to the model:

\begin{Shaded}
\begin{Highlighting}[]
\CommentTok{\# + absdiff("grade")}
\end{Highlighting}
\end{Shaded}

This could be computed like this:
\(\text{absdiff}(\text{"grade"}) = \sum_{i \neq j} y_{ij} \cdot |x_i^{\text{grade}} - x_j^{\text{grade}}|\),
where \(y_{ij}=1\) if student \(i\) has nominated student \(j\) and is
\(0\) otherwise. \(x_k^{\text{grade}}\) corresponds to the grade of
student \(k\). A negative coefficient would support the hypothesis in
this case, as it would make ties more probable if they have similar
grade levels. In the network graph, this effect would be seen as more
ties existing between students of similar academic level. It should also
be noted, that since we just assume one total grade for each student,
there might be different results if the same hypothesis is executed on
just single subject grades.

The second hypothesis is that students with higher grades receive more
friendship nominations. That would require this addition to the model to
be figured out if the hypothesis is supported or not:

\begin{Shaded}
\begin{Highlighting}[]
\CommentTok{\# + nodeicov("grade")}
\end{Highlighting}
\end{Shaded}

This could be computed like this:
\(\text{nodeicov}(\text{"grade"}) = \sum_{i \neq j} y_{ij} \cdot x_j^{\text{grade}}\).
A positive coefficient would support this hypothesis, as that would mean
that students with higher grades receive more nominations. In the
network graph, this could be seen as nodes with a high grade value
having a relatively higher indegree than those with a lower grade.

Both of these hypotheses could explain some residual variation in tie
formation, that has not yet been captured by our current model.

\section{Task 3: Estimation and interpretation of
SAOMs}\label{task-3-estimation-and-interpretation-of-saoms}

\subsection{3.1}\label{section-9}

\subsubsection{3.1.1}\label{section-10}

\begin{Shaded}
\begin{Highlighting}[]
\FunctionTok{install.packages}\NormalTok{(}\StringTok{"RSiena"}\NormalTok{)}
\end{Highlighting}
\end{Shaded}

\begin{Shaded}
\begin{Highlighting}[]
\FunctionTok{library}\NormalTok{(RSiena)}
\end{Highlighting}
\end{Shaded}

\begin{verbatim}
## Registered S3 method overwritten by 'RSiena':
##   method        from     
##   plot.sienaGOF autograph
\end{verbatim}

\begin{verbatim}
## 
## Attache Paket: 'RSiena'
\end{verbatim}

\begin{verbatim}
## Das folgende Objekt ist maskiert 'package:xtable':
## 
##     xtable
\end{verbatim}

\begin{Shaded}
\begin{Highlighting}[]
\FunctionTok{library}\NormalTok{(parallel)}
\FunctionTok{library}\NormalTok{(sna)}

\CommentTok{\# Load friendship waves}
\NormalTok{f1 }\OtherTok{\textless{}{-}} \FunctionTok{matrix}\NormalTok{(}\FunctionTok{as.numeric}\NormalTok{(}\FunctionTok{as.matrix}\NormalTok{(}\FunctionTok{read\_matrix}\NormalTok{(}\StringTok{"Glasgow/f1.csv"}\NormalTok{, }\AttributeTok{header =} \ConstantTok{FALSE}\NormalTok{))), }
             \AttributeTok{nrow =} \DecValTok{129}\NormalTok{, }\AttributeTok{ncol =} \DecValTok{129}\NormalTok{)}
\NormalTok{f2 }\OtherTok{\textless{}{-}} \FunctionTok{matrix}\NormalTok{(}\FunctionTok{as.numeric}\NormalTok{(}\FunctionTok{as.matrix}\NormalTok{(}\FunctionTok{read\_matrix}\NormalTok{(}\StringTok{"Glasgow/f2.csv"}\NormalTok{, }\AttributeTok{header =} \ConstantTok{FALSE}\NormalTok{))), }
             \AttributeTok{nrow =} \DecValTok{129}\NormalTok{, }\AttributeTok{ncol =} \DecValTok{129}\NormalTok{)}
\NormalTok{f3 }\OtherTok{\textless{}{-}} \FunctionTok{matrix}\NormalTok{(}\FunctionTok{as.numeric}\NormalTok{(}\FunctionTok{as.matrix}\NormalTok{(}\FunctionTok{read\_matrix}\NormalTok{(}\StringTok{"Glasgow/f3.csv"}\NormalTok{, }\AttributeTok{header =} \ConstantTok{FALSE}\NormalTok{))), }
             \AttributeTok{nrow =} \DecValTok{129}\NormalTok{, }\AttributeTok{ncol =} \DecValTok{129}\NormalTok{)}

\CommentTok{\# Read demographic, alcohol and logdistance data}
\NormalTok{demographic }\OtherTok{\textless{}{-}} \FunctionTok{read.csv}\NormalTok{(}\StringTok{"Glasgow/demographic.csv"}\NormalTok{, }\AttributeTok{header =} \ConstantTok{TRUE}\NormalTok{)}
\NormalTok{alcohol\_data }\OtherTok{\textless{}{-}} \FunctionTok{read.csv}\NormalTok{(}\StringTok{"Glasgow/alcohol.csv"}\NormalTok{, }\AttributeTok{header =} \ConstantTok{TRUE}\NormalTok{)}
\NormalTok{logdist }\OtherTok{\textless{}{-}} \FunctionTok{matrix}\NormalTok{(}\FunctionTok{as.numeric}\NormalTok{(}\FunctionTok{unlist}\NormalTok{(}\FunctionTok{strsplit}\NormalTok{(}\FunctionTok{readLines}\NormalTok{(}\StringTok{"Glasgow/logdistance.csv"}\NormalTok{), }\StringTok{","}\NormalTok{))),}\AttributeTok{nrow =} \DecValTok{129}\NormalTok{, }\AttributeTok{ncol =} \DecValTok{129}\NormalTok{)}


\NormalTok{friendship }\OtherTok{\textless{}{-}} \FunctionTok{as\_dependent\_rsiena}\NormalTok{(}\FunctionTok{array}\NormalTok{(}\FunctionTok{c}\NormalTok{(f1, f2, f3), }\AttributeTok{dim =} \FunctionTok{c}\NormalTok{(}\DecValTok{129}\NormalTok{, }\DecValTok{129}\NormalTok{, }\DecValTok{3}\NormalTok{)), }\AttributeTok{allowOnly =} \ConstantTok{FALSE}\NormalTok{)}

\CommentTok{\# Create covariates}
\NormalTok{gender }\OtherTok{\textless{}{-}} \FunctionTok{as\_covariate\_rsiena}\NormalTok{(demographic}\SpecialCharTok{$}\NormalTok{gender, }\AttributeTok{type =} \StringTok{"monadic"}\NormalTok{)}
\NormalTok{alcohol }\OtherTok{\textless{}{-}} \FunctionTok{as\_covariate\_rsiena}\NormalTok{(alcohol\_data[, }\DecValTok{1}\NormalTok{], }\AttributeTok{type =} \StringTok{"monadic"}\NormalTok{)}
\NormalTok{logdistance }\OtherTok{\textless{}{-}} \FunctionTok{as\_covariate\_rsiena}\NormalTok{(logdist, }\AttributeTok{type =} \StringTok{"oneMode"}\NormalTok{)}

\CommentTok{\# Combine into siena data object}
\NormalTok{mydata }\OtherTok{\textless{}{-}} \FunctionTok{make\_data\_rsiena}\NormalTok{(friendship, gender, alcohol, logdistance)}
\CommentTok{\#print(mydata)}
\end{Highlighting}
\end{Shaded}

\begin{Shaded}
\begin{Highlighting}[]
\CommentTok{\# Compute Jaccard index and data description}
\FunctionTok{write\_report}\NormalTok{(mydata, }\AttributeTok{outputName =} \StringTok{"glasgowJaccardIndex"}\NormalTok{)}
\end{Highlighting}
\end{Shaded}

From the data generated in the file ``glasgowJaccardIndex'', we have the
resulting Jaccard Index values from Wave 1 \(\rightarrow\) Wave 2 with a
value of \(0.304\) and from Wave 2 \(\rightarrow\) Wave 3 with a value
of \(0.351\). Both of these values fit neatly into the recommended range
of \([0.2,0.6]\), indicating that the data contains sufficient data to
estimate a SAOM. If the value were below that range, it would mean the
network changes too drastically and leaves no network stability, while a
value above the range would indicate that the network hardly changes at
all and shows no evolutionary traits. The higher Jaccard index from Wave
2 \(\rightarrow\) Wave 3 indicates that the network gained stability
over time.

\subsubsection{3.1.2}\label{section-11}

\begin{Shaded}
\begin{Highlighting}[]
\NormalTok{myeff }\OtherTok{\textless{}{-}} \FunctionTok{make\_specification}\NormalTok{(mydata)}

\CommentTok{\#endogenous structural effects, from what generally should hold for networks}
\NormalTok{myeff }\OtherTok{\textless{}{-}} \FunctionTok{set\_effect}\NormalTok{(myeff, }\FunctionTok{list}\NormalTok{(recip))       }\CommentTok{\# reciprocity}
\end{Highlighting}
\end{Shaded}

\begin{verbatim}
##   effectNumber effectName  shortName include fix   test  initialValue parm
## 1 17           reciprocity recip     TRUE    FALSE FALSE          0   0
\end{verbatim}

\begin{Shaded}
\begin{Highlighting}[]
\NormalTok{myeff }\OtherTok{\textless{}{-}} \FunctionTok{set\_effect}\NormalTok{(myeff, }\FunctionTok{list}\NormalTok{(transTrip))   }\CommentTok{\# transitivity}
\end{Highlighting}
\end{Shaded}

\begin{verbatim}
##   effectNumber effectName          shortName include fix   test  initialValue
## 1 23           transitive triplets transTrip TRUE    FALSE FALSE          0  
##   parm
## 1 0
\end{verbatim}

\begin{Shaded}
\begin{Highlighting}[]
\NormalTok{myeff }\OtherTok{\textless{}{-}} \FunctionTok{set\_effect}\NormalTok{(myeff, }\FunctionTok{list}\NormalTok{(cycle3))      }\CommentTok{\# 3{-}cycles}
\end{Highlighting}
\end{Shaded}

\begin{verbatim}
##   effectNumber effectName shortName include fix   test  initialValue parm
## 1 43           3-cycles   cycle3    TRUE    FALSE FALSE          0   0
\end{verbatim}

\begin{Shaded}
\begin{Highlighting}[]
\CommentTok{\#exogenous control effects, we know from task 2 that gender plays a huge role}
\NormalTok{myeff }\OtherTok{\textless{}{-}} \FunctionTok{set\_effect}\NormalTok{(myeff, }\FunctionTok{list}\NormalTok{(egoX, altX, sameX), }\AttributeTok{covar1 =} \StringTok{"gender"}\NormalTok{)}
\end{Highlighting}
\end{Shaded}

\begin{verbatim}
##   effectNumber effectName   shortName include fix   test  initialValue parm
## 1 297          gender alter altX      TRUE    FALSE FALSE          0   0   
## 2 312          gender ego   egoX      TRUE    FALSE FALSE          0   0   
## 3 368          same gender  sameX     TRUE    FALSE FALSE          0   0
\end{verbatim}

\begin{Shaded}
\begin{Highlighting}[]
\CommentTok{\#Hypothesis i: popularity}
\NormalTok{myeff }\OtherTok{\textless{}{-}} \FunctionTok{set\_effect}\NormalTok{(myeff, }\FunctionTok{list}\NormalTok{(inPopSqrt))}
\end{Highlighting}
\end{Shaded}

\begin{verbatim}
##   effectNumber effectName                   shortName include fix   test 
## 1 86           indegree - popularity (sqrt) inPopSqrt TRUE    FALSE FALSE
##   initialValue parm
## 1          0   0
\end{verbatim}

\begin{Shaded}
\begin{Highlighting}[]
\CommentTok{\#Hypothesis ii: alcohol homophily}
\NormalTok{myeff }\OtherTok{\textless{}{-}} \FunctionTok{set\_effect}\NormalTok{(myeff, }\FunctionTok{list}\NormalTok{(egoX, altX, simX), }\AttributeTok{covar1 =} \StringTok{"alcohol"}\NormalTok{)}
\end{Highlighting}
\end{Shaded}

\begin{verbatim}
##   effectNumber effectName         shortName include fix   test  initialValue
## 1 481          alcohol alter      altX      TRUE    FALSE FALSE          0  
## 2 496          alcohol ego        egoX      TRUE    FALSE FALSE          0  
## 3 542          alcohol similarity simX      TRUE    FALSE FALSE          0  
##   parm
## 1 0   
## 2 0   
## 3 0
\end{verbatim}

\begin{Shaded}
\begin{Highlighting}[]
\CommentTok{\#Hypothesis iii: geographic proximity}
\NormalTok{myeff }\OtherTok{\textless{}{-}} \FunctionTok{set\_effect}\NormalTok{(myeff, }\FunctionTok{list}\NormalTok{(X), }\AttributeTok{covar1 =} \StringTok{"logdistance"}\NormalTok{)}
\end{Highlighting}
\end{Shaded}

\begin{verbatim}
##   effectNumber effectName  shortName include fix   test  initialValue parm
## 1 264          logdistance X         TRUE    FALSE FALSE          0   0
\end{verbatim}

\begin{Shaded}
\begin{Highlighting}[]
\NormalTok{myeff}
\end{Highlighting}
\end{Shaded}

\begin{verbatim}
##    effectName                          include fix   test  initialValue parm
## 1  constant friendship rate (period 1) TRUE    FALSE FALSE    7.45424   0   
## 2  constant friendship rate (period 2) TRUE    FALSE FALSE    6.82927   0   
## 3  outdegree (density)                 TRUE    FALSE FALSE   -1.61299   0   
## 4  reciprocity                         TRUE    FALSE FALSE    0.00000   0   
## 5  transitive triplets                 TRUE    FALSE FALSE    0.00000   0   
## 6  3-cycles                            TRUE    FALSE FALSE    0.00000   0   
## 7  indegree - popularity (sqrt)        TRUE    FALSE FALSE    0.00000   0   
## 8  logdistance                         TRUE    FALSE FALSE    0.00000   0   
## 9  gender alter                        TRUE    FALSE FALSE    0.00000   0   
## 10 gender ego                          TRUE    FALSE FALSE    0.00000   0   
## 11 same gender                         TRUE    FALSE FALSE    0.00000   0   
## 12 alcohol alter                       TRUE    FALSE FALSE    0.00000   0   
## 13 alcohol ego                         TRUE    FALSE FALSE    0.00000   0   
## 14 alcohol similarity                  TRUE    FALSE FALSE    0.00000   0
\end{verbatim}

\subsubsection{3.1.3}\label{section-12}

\begin{Shaded}
\begin{Highlighting}[]
\CommentTok{\#use different algorithm package thingy to allow constraint.}
\NormalTok{SAOMAlgorithm }\OtherTok{\textless{}{-}} \FunctionTok{sienaAlgorithmCreate}\NormalTok{(}\AttributeTok{seed =} \DecValTok{161}\NormalTok{, }\AttributeTok{n3 =} \DecValTok{3000}\NormalTok{, }\AttributeTok{nsub =} \DecValTok{4}\NormalTok{, }\AttributeTok{MaxDegree =} \FunctionTok{c}\NormalTok{(}\AttributeTok{friendship =} \DecValTok{6}\NormalTok{))}
\end{Highlighting}
\end{Shaded}

\begin{verbatim}
## If you use this algorithm object, siena07 will create/use an output file Siena.txt .
\end{verbatim}

\begin{Shaded}
\begin{Highlighting}[]
\CommentTok{\# Estimate the model}
\end{Highlighting}
\end{Shaded}

\begin{Shaded}
\begin{Highlighting}[]
\NormalTok{model\_saom }\OtherTok{\textless{}{-}} \FunctionTok{siena07}\NormalTok{(SAOMAlgorithm,}
                \AttributeTok{data =}\NormalTok{ mydata, }
                \AttributeTok{effects =}\NormalTok{ myeff, }
                \AttributeTok{returnDeps =} \ConstantTok{TRUE}\NormalTok{, }
                \AttributeTok{batch =} \ConstantTok{FALSE}\NormalTok{,}
\NormalTok{                )}
\FunctionTok{saveRDS}\NormalTok{(model\_saom, }\StringTok{"model\_saom.rds"}\NormalTok{)}
\end{Highlighting}
\end{Shaded}

\begin{Shaded}
\begin{Highlighting}[]
\NormalTok{model\_saom }\OtherTok{\textless{}{-}} \FunctionTok{readRDS}\NormalTok{(}\StringTok{"model\_saom.rds"}\NormalTok{)}
\CommentTok{\#Check results and convergence criterion (as presented in the lecture slides)}
\CommentTok{\#model\_saom}
\NormalTok{t\_ratios }\OtherTok{\textless{}{-}} \FunctionTok{apply}\NormalTok{(model\_saom}\SpecialCharTok{$}\NormalTok{sf, }\DecValTok{2}\NormalTok{, mean) }\SpecialCharTok{/} \FunctionTok{apply}\NormalTok{(model\_saom}\SpecialCharTok{$}\NormalTok{sf, }\DecValTok{2}\NormalTok{, sd)}
\NormalTok{t\_ratios}
\end{Highlighting}
\end{Shaded}

\begin{verbatim}
## [1] -0.048280819 -0.046630219 -0.044989376 -0.055399050 -0.000493128
## [6]  0.036705480 -0.047920270 -0.016010575
\end{verbatim}

\begin{Shaded}
\begin{Highlighting}[]
\NormalTok{overall }\OtherTok{\textless{}{-}} \FunctionTok{sqrt}\NormalTok{(}\FunctionTok{t}\NormalTok{(}\FunctionTok{apply}\NormalTok{(model\_saom}\SpecialCharTok{$}\NormalTok{sf, }\DecValTok{2}\NormalTok{, mean)) }\SpecialCharTok{\%*\%} \FunctionTok{solve}\NormalTok{(}\FunctionTok{cov}\NormalTok{(model\_saom}\SpecialCharTok{$}\NormalTok{sf)) }\SpecialCharTok{\%*\%} \FunctionTok{apply}\NormalTok{(model\_saom}\SpecialCharTok{$}\NormalTok{sf, }\DecValTok{2}\NormalTok{, mean))}
\NormalTok{overall}
\end{Highlighting}
\end{Shaded}

\begin{verbatim}
##           [,1]
## [1,] 0.1158141
\end{verbatim}

As we can see from these results, we have that all absolute values of
the t-ratios are below \(0.1\), along with the overall value being below
\(0.25\). This fulfills the convergence criterion, indicating that all
of our chosen parameters were sufficient to achieve convergence. The
SAOM was estimated with 4 subphases and 3000 iterations in phase 3,
which was enough to achieve this convergence and did not need to have
another re-estimation or other model adjustments. These parameters
consist of the three needed for answering our hypotheses: One parameter
to control for the popularity of a student by checking for its indegree
value, another to check for alcohol consumption homophily and the last
one te see if geographic proximity has an affect on the model evolution.
Furthermore we have some basic endogenous and exogenous variables that
normally play a big part in the structure of most networks: Variables to
check for the reciprocity of ties, the transitivity of ties and 3-cyclic
patterns. Lastly we added in one more parameter, to give control for
gender homophily. We have seen in task 2 that it plays a huge role in
the structure of this specific network.

\subsubsection{3.1.4}\label{section-13}

\begin{Shaded}
\begin{Highlighting}[]
\CommentTok{\#most of this is just from the practical:}
\NormalTok{cl }\OtherTok{\textless{}{-}} \FunctionTok{makeCluster}\NormalTok{(}\DecValTok{4}\NormalTok{)}

\CommentTok{\# Indegree distribution}
\NormalTok{gof\_id }\OtherTok{\textless{}{-}} \FunctionTok{test\_gof}\NormalTok{(}
\NormalTok{    model\_saom,}
    \AttributeTok{verbose =} \ConstantTok{FALSE}\NormalTok{,}
    \AttributeTok{varName =} \StringTok{"friendship"}\NormalTok{, IndegreeDistribution,}
    \AttributeTok{cluster =}\NormalTok{ cl}
\NormalTok{)}

\CommentTok{\# Outdegree distribution}
\NormalTok{gof\_od }\OtherTok{\textless{}{-}} \FunctionTok{test\_gof}\NormalTok{(}
\NormalTok{    model\_saom,}
    \AttributeTok{verbose =} \ConstantTok{FALSE}\NormalTok{,}
    \AttributeTok{varName =} \StringTok{"friendship"}\NormalTok{, OutdegreeDistribution,}
    \AttributeTok{cluster =}\NormalTok{ cl}
\NormalTok{)}

\CommentTok{\# Triad census}
\NormalTok{gof\_tc }\OtherTok{\textless{}{-}} \FunctionTok{test\_gof}\NormalTok{(}
\NormalTok{    model\_saom,}
    \AttributeTok{verbose =} \ConstantTok{FALSE}\NormalTok{,}
    \AttributeTok{varName =} \StringTok{"friendship"}\NormalTok{, TriadCensus,}
    \AttributeTok{cluster =}\NormalTok{ cl}
\NormalTok{)}

\CommentTok{\# Geodesic distance}
\NormalTok{GeodesicDistribution }\OtherTok{\textless{}{-}} \ControlFlowTok{function}\NormalTok{(}
\NormalTok{    i, data, sims, period, groupName,}
\NormalTok{    varName, }\AttributeTok{levls =} \FunctionTok{c}\NormalTok{(}\DecValTok{1}\SpecialCharTok{:}\DecValTok{5}\NormalTok{, }\ConstantTok{Inf}\NormalTok{), }\AttributeTok{cumulative =} \ConstantTok{TRUE}
\NormalTok{) \{}
\NormalTok{    x }\OtherTok{\textless{}{-}} \FunctionTok{networkExtraction}\NormalTok{(i, data, sims, period, groupName, varName)}
    \FunctionTok{require}\NormalTok{(sna)}
\NormalTok{    a }\OtherTok{\textless{}{-}}\NormalTok{ sna}\SpecialCharTok{::}\FunctionTok{geodist}\NormalTok{(}\FunctionTok{symmetrize}\NormalTok{(x))}\SpecialCharTok{$}\NormalTok{gdist}
    \ControlFlowTok{if}\NormalTok{ (cumulative) \{}
\NormalTok{        gdi }\OtherTok{\textless{}{-}} \FunctionTok{sapply}\NormalTok{(levls, }\ControlFlowTok{function}\NormalTok{(i) \{ }\FunctionTok{sum}\NormalTok{(a }\SpecialCharTok{\textless{}=}\NormalTok{ i) \})}
\NormalTok{    \} }\ControlFlowTok{else}\NormalTok{ \{}
\NormalTok{        gdi }\OtherTok{\textless{}{-}} \FunctionTok{sapply}\NormalTok{(levls, }\ControlFlowTok{function}\NormalTok{(i) \{ }\FunctionTok{sum}\NormalTok{(a }\SpecialCharTok{==}\NormalTok{ i) \})}
\NormalTok{    \}}
    \FunctionTok{names}\NormalTok{(gdi) }\OtherTok{\textless{}{-}} \FunctionTok{as.character}\NormalTok{(levls)}
\NormalTok{    gdi}
\NormalTok{\}}

\NormalTok{gof\_gd }\OtherTok{\textless{}{-}} \FunctionTok{test\_gof}\NormalTok{(}
\NormalTok{    model\_saom,}
    \AttributeTok{verbose =} \ConstantTok{FALSE}\NormalTok{,}
    \AttributeTok{varName =} \StringTok{"friendship"}\NormalTok{, GeodesicDistribution}
\NormalTok{)}

\FunctionTok{stopCluster}\NormalTok{(cl)}

\CommentTok{\#Plot all four}
\FunctionTok{plot}\NormalTok{(gof\_id)}
\end{Highlighting}
\end{Shaded}

\pandocbounded{\includegraphics[keepaspectratio]{Submission2_files/Submission2_files/figure-latex/unnamed-chunk-25-1.pdf}}

\begin{Shaded}
\begin{Highlighting}[]
\FunctionTok{plot}\NormalTok{(gof\_od)}
\end{Highlighting}
\end{Shaded}

\begin{verbatim}
## Note: some statistics are not plotted because their variance is 0.
## This holds for the statistics: 6 7 8.
\end{verbatim}

\pandocbounded{\includegraphics[keepaspectratio]{Submission2_files/Submission2_files/figure-latex/unnamed-chunk-25-2.pdf}}

\begin{Shaded}
\begin{Highlighting}[]
\FunctionTok{plot}\NormalTok{(gof\_gd)}
\end{Highlighting}
\end{Shaded}

\begin{verbatim}
## Note: some statistics are not plotted because their variance is 0.
## This holds for the statistic: Inf.
\end{verbatim}

\pandocbounded{\includegraphics[keepaspectratio]{Submission2_files/Submission2_files/figure-latex/unnamed-chunk-25-3.pdf}}

\begin{Shaded}
\begin{Highlighting}[]
\FunctionTok{plot}\NormalTok{(gof\_tc, }\AttributeTok{center =} \ConstantTok{TRUE}\NormalTok{, }\AttributeTok{scale =} \ConstantTok{TRUE}\NormalTok{)}
\end{Highlighting}
\end{Shaded}

\pandocbounded{\includegraphics[keepaspectratio]{Submission2_files/Submission2_files/figure-latex/unnamed-chunk-25-4.pdf}}

The plot for the goodness of fit for the indegree distribution shows
that our SAOM model captures this property well. The plot follows the
real data and as such produces a p-value of \(0.974\), which supports
that the model captures this property well, since high p-values close to
\(1\) support the property being captured. The Outdegree distribution is
at a low p-value of \(0.047\), which means it somewhat captures this
property. This had a p-value of 0 before implementing the maximum
outdegree constraint of \(6\) per student, as the data collection method
enforces, meaning that the constraint is very important to capture the
outdegree structure of the observed network, but it seems there are
still some other variables that are important for the developement of
the network that were not accounted for so far.

The other two plotted parameters are both not captured well. To start of
with, they both have a p-value of \(0\), which means that neither of
them accurately capture the observed data. The geodesic distribution
itself has values that are higher than they are supposed to be over all
distances, granting it a p-value of \(0\). This means that the simulated
network is more connected than the real one, and that the real one has
more of a community type structured one, meaning there are more seperate
clusters that have only few edges going out / coming in. The Triad
consensus also has a p-value of \(0\), and as seen in the graph, while
it correctly predicts some kinds of triads, it is completely wrong by
overestimating some and underestimating others. It makes it clear that
some important terms for the development of the network are still left
out and not being considered.

\subsubsection{3.1.5}\label{section-14}

\begin{Shaded}
\begin{Highlighting}[]
\FunctionTok{source}\NormalTok{(}\StringTok{"printSiena.R"}\NormalTok{)}
\NormalTok{model\_saom}
\end{Highlighting}
\end{Shaded}

\begin{verbatim}
## Estimates, standard errors and convergence t-ratios
## 
##                                         Estimate   Standard   Convergence 
##                                                      Error      t-ratio   
## 
## Rate parameters: 
##   0.1      Rate parameter period 1      12.3619  ( 1.3812   )             
##   0.2      Rate parameter period 2       9.7971  ( 0.9900   )             
## 
## Other parameters: 
##   1.  eval outdegree (density)          -2.5681  ( 0.1455   )   -0.0483   
##   2.  eval reciprocity                   2.2687  ( 0.0867   )   -0.0466   
##   3.  eval 3-cycles                      1.0295  ( 0.0537   )   -0.0450   
##   4.  eval indegree - popularity (sqrt)  0.0056  ( 0.0689   )   -0.0554   
##   5.  eval logdistance                  -0.1852  ( 0.0447   )   -0.0005   
##   6.  eval alcohol alter                 0.0017  ( 0.0314   )    0.0367   
##   7.  eval alcohol ego                  -0.0561  ( 0.0385   )   -0.0479   
##   8.  eval alcohol similarity            0.4939  ( 0.1650   )   -0.0160   
## 
## Overall maximum convergence ratio:    0.1158 
## 
## 
## Degrees constrained to maximum values:
## friendship : 6 
## 
## 
## Total of 4510 iteration steps.
\end{verbatim}

\begin{Shaded}
\begin{Highlighting}[]
\FunctionTok{printSiena}\NormalTok{(model\_saom)}
\end{Highlighting}
\end{Shaded}

\begin{verbatim}
##     dependent                              effect  theta  s.e. p.value sig.
## 1        rate constant friendship rate (period 1) 12.362 1.381             
## 2        rate constant friendship rate (period 2)  9.797 0.990             
## 3  friendship                 outdegree (density) -2.568 0.145       0  ***
## 4  friendship                         reciprocity  2.269 0.087       0  ***
## 5  friendship                            3-cycles  1.029 0.054       0  ***
## 6  friendship        indegree - popularity (sqrt)  0.006 0.069   0.935     
## 7  friendship                         logdistance -0.185 0.045       0  ***
## 8  friendship                       alcohol alter  0.002 0.031   0.956     
## 9  friendship                         alcohol ego -0.056 0.039   0.145     
## 10 friendship                  alcohol similarity  0.494 0.165   0.003   **
##    t.conv
## 1        
## 2        
## 3  -0.048
## 4  -0.047
## 5  -0.045
## 6  -0.055
## 7       0
## 8   0.037
## 9  -0.048
## 10 -0.016
\end{verbatim}

Hypothesis (\(i\)) Students tend to be friends with popular pupils: This
hypothesis is not supported by the SAOM model. The extremely high
p-value \(0.935\) shows no significance, along with the the near zero
coefficient \(+0.006\), showing that even if it had significance, the
effect would be barely noticeable. Since SAOM estimate how the network
changes, it just says the popularity plays no role in the formation of
new ties. A student will nominate a person regardless of how popular
that student already is.

Hypothesis (\(ii\)) Students tend to be friends with pupils with similar
alcohol consumption to theirs: With a p-value of \(0.003\), this
hypothesis is significant and with a coefficient value of \(+0.494\) is
supported. The likelyhood of a tie forming between two students at the
same alcohol consumption level is \(e^{0.494} \approx 1.64\) times
higher than if they had the biggest possible difference. This means that
alcohol homophily is an important factor in the development of the
structural network and that friendships will center partly around the
alcohol consumption level of students. This means that alcohol
consumption contributes to cluster formation around alcohol homophily.
It is to be noted that both alcohol alter and alcohol ego have high
p-values and low total coefficients, meaning that alcohol consumption
has no noticeable effect on the amount of nominations a student will
send or receive.

Hypothesis (\(iii\)) Students tend to be friends with students who live
in the same neighborhood (living nearby): The values of logdistance
support this hypothesis. With a near zero p-value and a coefficient of
\(-0.185\), the simulated model predicts that the further away two
students live from each other, the less likely it becomes that they
nominate one another. Thus they are more likely to nominate someone who
lives close to them. Similar to alcohol homophily, this means that the
geological position of students homes will lead to some clustering over
enough time. So if the network has its nodes centered on a map, it
should show clear clustering around certain areas like different
neighborhoods.

\subsection{3.2}\label{section-15}

\subsubsection{3.2.1}\label{section-16}

\begin{Shaded}
\begin{Highlighting}[]
\CommentTok{\# Alcohol as dependent behaviour variable}
\NormalTok{alcoholbeh }\OtherTok{\textless{}{-}} \FunctionTok{as\_dependent\_rsiena}\NormalTok{(}
  \FunctionTok{as.matrix}\NormalTok{(alcohol\_data[, }\DecValTok{1}\SpecialCharTok{:}\DecValTok{3}\NormalTok{]), }
  \AttributeTok{type =} \StringTok{"behavior"}
\NormalTok{)}

\CommentTok{\# Rebuild data object with alcohol as dependent variable}
\NormalTok{mydata\_coev }\OtherTok{\textless{}{-}} \FunctionTok{make\_data\_rsiena}\NormalTok{(friendship, alcoholbeh, gender, logdistance)}

\CommentTok{\# Check preconditions}
\end{Highlighting}
\end{Shaded}

\begin{Shaded}
\begin{Highlighting}[]
\FunctionTok{write\_report}\NormalTok{(mydata\_coev, }\AttributeTok{outputName =} \StringTok{"glasgowCoevReport"}\NormalTok{)}
\end{Highlighting}
\end{Shaded}

\begin{Shaded}
\begin{Highlighting}[]
\CommentTok{\# Moran\textquotesingle{}s I for autocorrelation in alcohol behaviour}
\NormalTok{moran1 }\OtherTok{\textless{}{-}} \FunctionTok{nacf}\NormalTok{(f1, alcohol\_data[, }\DecValTok{1}\NormalTok{], }\AttributeTok{lag.max =} \DecValTok{1}\NormalTok{, }\AttributeTok{type =} \StringTok{"moran"}\NormalTok{,}
               \AttributeTok{neighborhood.type =} \StringTok{"out"}\NormalTok{, }\AttributeTok{mode =} \StringTok{"digraph"}\NormalTok{)}
\NormalTok{moran2 }\OtherTok{\textless{}{-}} \FunctionTok{nacf}\NormalTok{(f2, alcohol\_data[, }\DecValTok{2}\NormalTok{], }\AttributeTok{lag.max =} \DecValTok{1}\NormalTok{, }\AttributeTok{type =} \StringTok{"moran"}\NormalTok{,}
               \AttributeTok{neighborhood.type =} \StringTok{"out"}\NormalTok{, }\AttributeTok{mode =} \StringTok{"digraph"}\NormalTok{)}
\NormalTok{moran3 }\OtherTok{\textless{}{-}} \FunctionTok{nacf}\NormalTok{(f3, alcohol\_data[, }\DecValTok{3}\NormalTok{], }\AttributeTok{lag.max =} \DecValTok{1}\NormalTok{, }\AttributeTok{type =} \StringTok{"moran"}\NormalTok{,}
               \AttributeTok{neighborhood.type =} \StringTok{"out"}\NormalTok{, }\AttributeTok{mode =} \StringTok{"digraph"}\NormalTok{)}
\NormalTok{autocorr }\OtherTok{\textless{}{-}} \FunctionTok{rbind}\NormalTok{(moran1, moran2, moran3)}
\NormalTok{autocorr[, }\DecValTok{2}\NormalTok{]}
\end{Highlighting}
\end{Shaded}

\begin{verbatim}
##    moran1    moran2    moran3 
## 0.2450651 0.3394460 0.3316509
\end{verbatim}

\begin{Shaded}
\begin{Highlighting}[]
\DocumentationTok{\#\# Model specification}
\NormalTok{myeff\_coev }\OtherTok{\textless{}{-}} \FunctionTok{make\_specification}\NormalTok{(mydata\_coev)}

\CommentTok{\# {-}{-}{-} Selection part: identical to 3.1.2 {-}{-}{-}}
\NormalTok{myeff\_coev }\OtherTok{\textless{}{-}} \FunctionTok{set\_effect}\NormalTok{(myeff\_coev, }\FunctionTok{list}\NormalTok{(recip),       }\AttributeTok{depvar =} \StringTok{"friendship"}\NormalTok{)}
\end{Highlighting}
\end{Shaded}

\begin{verbatim}
##   effectNumber effectName  shortName include fix   test  initialValue parm
## 1 17           reciprocity recip     TRUE    FALSE FALSE          0   0
\end{verbatim}

\begin{Shaded}
\begin{Highlighting}[]
\NormalTok{myeff\_coev }\OtherTok{\textless{}{-}} \FunctionTok{set\_effect}\NormalTok{(myeff\_coev, }\FunctionTok{list}\NormalTok{(transTrip),   }\AttributeTok{depvar =} \StringTok{"friendship"}\NormalTok{)}
\end{Highlighting}
\end{Shaded}

\begin{verbatim}
##   effectNumber effectName          shortName include fix   test  initialValue
## 1 23           transitive triplets transTrip TRUE    FALSE FALSE          0  
##   parm
## 1 0
\end{verbatim}

\begin{Shaded}
\begin{Highlighting}[]
\NormalTok{myeff\_coev }\OtherTok{\textless{}{-}} \FunctionTok{set\_effect}\NormalTok{(myeff\_coev, }\FunctionTok{list}\NormalTok{(cycle3),      }\AttributeTok{depvar =} \StringTok{"friendship"}\NormalTok{)}
\end{Highlighting}
\end{Shaded}

\begin{verbatim}
##   effectNumber effectName shortName include fix   test  initialValue parm
## 1 43           3-cycles   cycle3    TRUE    FALSE FALSE          0   0
\end{verbatim}

\begin{Shaded}
\begin{Highlighting}[]
\NormalTok{myeff\_coev }\OtherTok{\textless{}{-}} \FunctionTok{set\_effect}\NormalTok{(myeff\_coev, }\FunctionTok{list}\NormalTok{(egoX, altX, sameX), }\AttributeTok{depvar =} \StringTok{"friendship"}\NormalTok{,}
                         \AttributeTok{covar1 =} \StringTok{"gender"}\NormalTok{)}
\end{Highlighting}
\end{Shaded}

\begin{verbatim}
##   effectNumber effectName   shortName include fix   test  initialValue parm
## 1 297          gender alter altX      TRUE    FALSE FALSE          0   0   
## 2 312          gender ego   egoX      TRUE    FALSE FALSE          0   0   
## 3 368          same gender  sameX     TRUE    FALSE FALSE          0   0
\end{verbatim}

\begin{Shaded}
\begin{Highlighting}[]
\CommentTok{\# Hypothesis i: popularity}
\NormalTok{myeff\_coev }\OtherTok{\textless{}{-}} \FunctionTok{set\_effect}\NormalTok{(myeff\_coev, }\FunctionTok{list}\NormalTok{(inPopSqrt),   }\AttributeTok{depvar =} \StringTok{"friendship"}\NormalTok{)}
\end{Highlighting}
\end{Shaded}

\begin{verbatim}
##   effectNumber effectName                   shortName include fix   test 
## 1 86           indegree - popularity (sqrt) inPopSqrt TRUE    FALSE FALSE
##   initialValue parm
## 1          0   0
\end{verbatim}

\begin{Shaded}
\begin{Highlighting}[]
\CommentTok{\# Hypothesis ii: alcohol homophily (now using the behaviour dependent variable)}
\NormalTok{myeff\_coev }\OtherTok{\textless{}{-}} \FunctionTok{set\_effect}\NormalTok{(myeff\_coev, }\FunctionTok{list}\NormalTok{(egoX, altX, simX), }\AttributeTok{depvar =} \StringTok{"friendship"}\NormalTok{,}
                         \AttributeTok{covar1 =} \StringTok{"alcoholbeh"}\NormalTok{)}
\end{Highlighting}
\end{Shaded}

\begin{verbatim}
##   effectNumber effectName            shortName include fix   test  initialValue
## 1 481          alcoholbeh alter      altX      TRUE    FALSE FALSE          0  
## 2 496          alcoholbeh ego        egoX      TRUE    FALSE FALSE          0  
## 3 542          alcoholbeh similarity simX      TRUE    FALSE FALSE          0  
##   parm
## 1 0   
## 2 0   
## 3 0
\end{verbatim}

\begin{Shaded}
\begin{Highlighting}[]
\CommentTok{\# Hypothesis iii: geographic proximity}
\NormalTok{myeff\_coev }\OtherTok{\textless{}{-}} \FunctionTok{set\_effect}\NormalTok{(myeff\_coev, }\FunctionTok{list}\NormalTok{(X), }\AttributeTok{depvar =} \StringTok{"friendship"}\NormalTok{,}
                         \AttributeTok{covar1 =} \StringTok{"logdistance"}\NormalTok{)}
\end{Highlighting}
\end{Shaded}

\begin{verbatim}
##   effectNumber effectName  shortName include fix   test  initialValue parm
## 1 264          logdistance X         TRUE    FALSE FALSE          0   0
\end{verbatim}

\begin{Shaded}
\begin{Highlighting}[]
\CommentTok{\# {-}{-}{-} Influence part: effects explaining alcohol behaviour evolution {-}{-}{-}}

\CommentTok{\# Basic controls for behaviour (linear and quadratic shape)}
\CommentTok{\# These are included by default in make\_specification but listed here for clarity}
\CommentTok{\# They capture the baseline tendency to increase/decrease alcohol consumption}

\CommentTok{\# Hypothesis iv: Popular students tend to increase or maintain alcohol consumption}
\CommentTok{\# indeg: students who receive more friendship nominations (more popular)}
\CommentTok{\# tend to increase their alcohol consumption}
\NormalTok{myeff\_coev }\OtherTok{\textless{}{-}} \FunctionTok{set\_effect}\NormalTok{(myeff\_coev, }\FunctionTok{list}\NormalTok{(indeg), }\AttributeTok{depvar =} \StringTok{"alcoholbeh"}\NormalTok{,}
                         \AttributeTok{covar1 =} \StringTok{"friendship"}\NormalTok{)}
\end{Highlighting}
\end{Shaded}

\begin{verbatim}
##   effectNumber effectName          shortName include fix   test  initialValue
## 1 792          alcoholbeh indegree indeg     TRUE    FALSE FALSE          0  
##   parm
## 1 0
\end{verbatim}

\begin{Shaded}
\begin{Highlighting}[]
\CommentTok{\# Hypothesis v: Students adjust alcohol consumption to that of their friends}
\CommentTok{\# avSim: average similarity effect — ego moves their behaviour toward}
\CommentTok{\# the average alcohol level of their alters (the core social influence effect)}
\NormalTok{myeff\_coev }\OtherTok{\textless{}{-}} \FunctionTok{set\_effect}\NormalTok{(myeff\_coev, }\FunctionTok{list}\NormalTok{(avSim), }\AttributeTok{depvar =} \StringTok{"alcoholbeh"}\NormalTok{,}
                         \AttributeTok{covar1 =} \StringTok{"friendship"}\NormalTok{)}
\end{Highlighting}
\end{Shaded}

\begin{verbatim}
##   effectNumber effectName                    shortName include fix   test 
## 1 755          alcoholbeh average similarity avSim     TRUE    FALSE FALSE
##   initialValue parm
## 1          0   0
\end{verbatim}

\begin{Shaded}
\begin{Highlighting}[]
\NormalTok{myeff\_coev}
\end{Highlighting}
\end{Shaded}

\begin{verbatim}
##    name       effectName                          include fix   test 
## 1  friendship constant friendship rate (period 1) TRUE    FALSE FALSE
## 2  friendship constant friendship rate (period 2) TRUE    FALSE FALSE
## 3  friendship outdegree (density)                 TRUE    FALSE FALSE
## 4  friendship reciprocity                         TRUE    FALSE FALSE
## 5  friendship transitive triplets                 TRUE    FALSE FALSE
## 6  friendship 3-cycles                            TRUE    FALSE FALSE
## 7  friendship indegree - popularity (sqrt)        TRUE    FALSE FALSE
## 8  friendship logdistance                         TRUE    FALSE FALSE
## 9  friendship gender alter                        TRUE    FALSE FALSE
## 10 friendship gender ego                          TRUE    FALSE FALSE
## 11 friendship same gender                         TRUE    FALSE FALSE
## 12 friendship alcoholbeh alter                    TRUE    FALSE FALSE
## 13 friendship alcoholbeh ego                      TRUE    FALSE FALSE
## 14 friendship alcoholbeh similarity               TRUE    FALSE FALSE
## 15 alcoholbeh rate alcoholbeh (period 1)          TRUE    FALSE FALSE
## 16 alcoholbeh rate alcoholbeh (period 2)          TRUE    FALSE FALSE
## 17 alcoholbeh alcoholbeh linear shape             TRUE    FALSE FALSE
## 18 alcoholbeh alcoholbeh quadratic shape          TRUE    FALSE FALSE
## 19 alcoholbeh alcoholbeh average similarity       TRUE    FALSE FALSE
## 20 alcoholbeh alcoholbeh indegree                 TRUE    FALSE FALSE
##    initialValue parm
## 1     7.45424   0   
## 2     6.82927   0   
## 3    -1.61299   0   
## 4     0.00000   0   
## 5     0.00000   0   
## 6     0.00000   0   
## 7     0.00000   0   
## 8     0.00000   0   
## 9     0.00000   0   
## 10    0.00000   0   
## 11    0.00000   0   
## 12    0.00000   0   
## 13    0.00000   0   
## 14    0.00000   0   
## 15    0.91000   0   
## 16    1.09314   0   
## 17    0.34704   0   
## 18    0.00000   0   
## 19    0.00000   0   
## 20    0.00000   0
\end{verbatim}

\subsubsection{3.2.2}\label{section-17}

\begin{Shaded}
\begin{Highlighting}[]
\CommentTok{\# first step model estimation}
\NormalTok{SAOMAlgorithm\_coev }\OtherTok{\textless{}{-}} \FunctionTok{sienaAlgorithmCreate}\NormalTok{(}\AttributeTok{seed =} \DecValTok{161}\NormalTok{, }\AttributeTok{n3 =} \DecValTok{3000}\NormalTok{, }\AttributeTok{nsub =} \DecValTok{4}\NormalTok{,}
                                           \AttributeTok{MaxDegree =} \FunctionTok{c}\NormalTok{(}\AttributeTok{friendship =} \DecValTok{6}\NormalTok{))}

\NormalTok{model\_coev }\OtherTok{\textless{}{-}} \FunctionTok{siena07}\NormalTok{(SAOMAlgorithm\_coev,}
                      \AttributeTok{data    =}\NormalTok{ mydata\_coev,}
                      \AttributeTok{effects =}\NormalTok{ myeff\_coev,}
                      \AttributeTok{returnDeps =} \ConstantTok{TRUE}\NormalTok{,}
                      \AttributeTok{batch =} \ConstantTok{FALSE}\NormalTok{)}

\CommentTok{\# second step model estimation}
\NormalTok{SAOMAlgorithm\_converge }\OtherTok{\textless{}{-}} \FunctionTok{sienaAlgorithmCreate}\NormalTok{(}
  \AttributeTok{seed =} \DecValTok{161}\NormalTok{, }
  \AttributeTok{n3 =} \DecValTok{5000}\NormalTok{, }
  \AttributeTok{nsub =} \DecValTok{5}\NormalTok{, }
  \AttributeTok{MaxDegree =} \FunctionTok{c}\NormalTok{(}\AttributeTok{friendship =} \DecValTok{6}\NormalTok{)}
\NormalTok{)}

\NormalTok{model\_coev\_final }\OtherTok{\textless{}{-}} \FunctionTok{siena07}\NormalTok{(}
\NormalTok{  SAOMAlgorithm\_converge,}
  \AttributeTok{data =}\NormalTok{ mydata\_coev,}
  \AttributeTok{effects =}\NormalTok{ myeff\_coev,}
  \AttributeTok{prevAns =}\NormalTok{ model\_coev, }\CommentTok{\# Using previous failed run as the starting point}
  \AttributeTok{returnDeps =} \ConstantTok{TRUE}\NormalTok{,}
  \AttributeTok{batch =} \ConstantTok{FALSE}
\NormalTok{)}

\CommentTok{\#model\_coev\_final}

\CommentTok{\# save model}
\FunctionTok{saveRDS}\NormalTok{(model\_coev\_final, }\StringTok{"model\_coef\_final.rds"}\NormalTok{)}
\end{Highlighting}
\end{Shaded}

\begin{Shaded}
\begin{Highlighting}[]
\NormalTok{model\_coev\_final }\OtherTok{\textless{}{-}} \FunctionTok{readRDS}\NormalTok{(}\StringTok{"model\_coef\_final.rds"}\NormalTok{)}
\CommentTok{\# Convergence check}
\NormalTok{t\_ratios\_coev }\OtherTok{\textless{}{-}} \FunctionTok{apply}\NormalTok{(model\_coev\_final}\SpecialCharTok{$}\NormalTok{sf, }\DecValTok{2}\NormalTok{, mean) }\SpecialCharTok{/} \FunctionTok{apply}\NormalTok{(model\_coev\_final}\SpecialCharTok{$}\NormalTok{sf, }\DecValTok{2}\NormalTok{, sd)}
\NormalTok{t\_ratios\_coev}
\end{Highlighting}
\end{Shaded}

\begin{verbatim}
##  [1] -0.0135874937 -0.0086024967  0.0300746609  0.0316140308  0.0165487528
##  [6]  0.0093000455  0.0339679526  0.0346385251 -0.0180897629 -0.0082636378
## [11]  0.0314324650  0.0079810074 -0.0037956706  0.0005432766  0.0269788564
## [16]  0.0168935341  0.0068595218 -0.0099553750 -0.0147690147  0.0081376447
\end{verbatim}

\begin{Shaded}
\begin{Highlighting}[]
\NormalTok{overall\_coev }\OtherTok{\textless{}{-}} \FunctionTok{sqrt}\NormalTok{(}\FunctionTok{t}\NormalTok{(}\FunctionTok{apply}\NormalTok{(model\_coev\_final}\SpecialCharTok{$}\NormalTok{sf, }\DecValTok{2}\NormalTok{, mean)) }\SpecialCharTok{\%*\%} 
                       \FunctionTok{solve}\NormalTok{(}\FunctionTok{cov}\NormalTok{(model\_coev\_final}\SpecialCharTok{$}\NormalTok{sf)) }\SpecialCharTok{\%*\%} 
                       \FunctionTok{apply}\NormalTok{(model\_coev\_final}\SpecialCharTok{$}\NormalTok{sf, }\DecValTok{2}\NormalTok{, mean))}
\NormalTok{overall\_coev}
\end{Highlighting}
\end{Shaded}

\begin{verbatim}
##            [,1]
## [1,] 0.09271135
\end{verbatim}

As can be seen from these results, all absolute values of the t-ratios
are below 0.1 and additionally the overall value also lies below 0.25.
This fulfills the convergence criterion, indicating that all of our
chosen parameters were sufficient to achieve convergence. The SAOM was
estimated with 5 subphases and 5000 iterations in phase 3, which was
enough to achieve this convergence and did not need to have another
re-estimation or other model adjustments. These parameters consist of
the three needed for answering our hypothesis: one parameter to control
the popularity of a student by checking for the indegree-popularity
value, another to check for alcohol consumption similarity and the last
one to see if geographic proximity has an effect on the model evolution.
For the influence part of the model we included average similarity and
indegree effects on the alcohol behavior dependent variable.
Furthermore, we included basic endogenous and exogenous variables that
play a significant part in the structure of most networks: variables to
check for reciprocity, transitive triplets, and 3-cyclic patterns. We
also added a parameter to control for same-gender homophily, as we have
seen that gender plays a major role in student network structures.

\subsubsection{3.2.3}\label{section-18}

\begin{Shaded}
\begin{Highlighting}[]
\DocumentationTok{\#\# Goodness of fit}
\NormalTok{cl }\OtherTok{\textless{}{-}} \FunctionTok{makeCluster}\NormalTok{(}\DecValTok{4}\NormalTok{)}

\NormalTok{gof\_coev\_id }\OtherTok{\textless{}{-}} \FunctionTok{test\_gof}\NormalTok{(model\_coev\_final, }\AttributeTok{verbose =} \ConstantTok{FALSE}\NormalTok{,}
                        \AttributeTok{varName =} \StringTok{"friendship"}\NormalTok{, IndegreeDistribution, }\AttributeTok{cluster =}\NormalTok{ cl)}
\NormalTok{gof\_coev\_od }\OtherTok{\textless{}{-}} \FunctionTok{test\_gof}\NormalTok{(model\_coev\_final, }\AttributeTok{verbose =} \ConstantTok{FALSE}\NormalTok{,}
                        \AttributeTok{varName =} \StringTok{"friendship"}\NormalTok{, OutdegreeDistribution, }\AttributeTok{cluster =}\NormalTok{ cl)}
\NormalTok{gof\_coev\_tc }\OtherTok{\textless{}{-}} \FunctionTok{test\_gof}\NormalTok{(model\_coev\_final, }\AttributeTok{verbose =} \ConstantTok{FALSE}\NormalTok{,}
                        \AttributeTok{varName =} \StringTok{"friendship"}\NormalTok{, TriadCensus, }\AttributeTok{cluster =}\NormalTok{ cl)}
\NormalTok{gof\_coev\_gd }\OtherTok{\textless{}{-}} \FunctionTok{test\_gof}\NormalTok{(model\_coev\_final, }\AttributeTok{verbose =} \ConstantTok{FALSE}\NormalTok{,}
                        \AttributeTok{varName =} \StringTok{"friendship"}\NormalTok{, GeodesicDistribution)}
\NormalTok{gof\_coev\_beh }\OtherTok{\textless{}{-}} \FunctionTok{test\_gof}\NormalTok{(model\_coev\_final, }\AttributeTok{verbose =} \ConstantTok{FALSE}\NormalTok{,}
                         \AttributeTok{varName =} \StringTok{"alcoholbeh"}\NormalTok{, BehaviorDistribution, }\AttributeTok{cluster =}\NormalTok{ cl)}

\FunctionTok{stopCluster}\NormalTok{(cl)}

\FunctionTok{plot}\NormalTok{(gof\_coev\_id)}
\end{Highlighting}
\end{Shaded}

\pandocbounded{\includegraphics[keepaspectratio]{Submission2_files/Submission2_files/figure-latex/unnamed-chunk-32-1.pdf}}

\begin{Shaded}
\begin{Highlighting}[]
\FunctionTok{plot}\NormalTok{(gof\_coev\_od)}
\end{Highlighting}
\end{Shaded}

\begin{verbatim}
## Note: some statistics are not plotted because their variance is 0.
## This holds for the statistics: 6 7 8.
\end{verbatim}

\pandocbounded{\includegraphics[keepaspectratio]{Submission2_files/Submission2_files/figure-latex/unnamed-chunk-32-2.pdf}}

\begin{Shaded}
\begin{Highlighting}[]
\FunctionTok{plot}\NormalTok{(gof\_coev\_gd)}
\end{Highlighting}
\end{Shaded}

\begin{verbatim}
## Note: some statistics are not plotted because their variance is 0.
## This holds for the statistic: Inf.
\end{verbatim}

\pandocbounded{\includegraphics[keepaspectratio]{Submission2_files/Submission2_files/figure-latex/unnamed-chunk-32-3.pdf}}

\begin{Shaded}
\begin{Highlighting}[]
\FunctionTok{plot}\NormalTok{(gof\_coev\_tc, }\AttributeTok{center =} \ConstantTok{TRUE}\NormalTok{, }\AttributeTok{scale =} \ConstantTok{TRUE}\NormalTok{)}
\end{Highlighting}
\end{Shaded}

\pandocbounded{\includegraphics[keepaspectratio]{Submission2_files/Submission2_files/figure-latex/unnamed-chunk-32-4.pdf}}

\begin{Shaded}
\begin{Highlighting}[]
\FunctionTok{plot}\NormalTok{(gof\_coev\_beh)}
\end{Highlighting}
\end{Shaded}

\begin{verbatim}
## Note: some statistics are not plotted because their variance is 0.
## This holds for the statistic: 5.
\end{verbatim}

\pandocbounded{\includegraphics[keepaspectratio]{Submission2_files/Submission2_files/figure-latex/unnamed-chunk-32-5.pdf}}

\begin{Shaded}
\begin{Highlighting}[]
\CommentTok{\#model\_coev\_final}
\end{Highlighting}
\end{Shaded}

The plot for the goodness of fit for the indegree distribution shows
that our SAOM model captures this property very well. The plot follows
the real data and as such produces a p-value of 0.933, which supports
the fact that the model captures this property well. This is because
high p-values close to 1 support the property being captured.
Additionally, the behavior distribution for alcohol consumption is
captured effectively with a p-value of 0.563, indicating the influence
part of the model is performing reasonably well. The Outdegree
distribution is at a low p-value of 0.01, which means it does not
capture this property well. This is despite implementing the maximum
outdegree constraint of 6 per student as the data collection method
enforces, meaning that while the constraint is very important to reflect
the observed network's limits, there are still some other variables
important for the development of the network that were not accounted for
so far. The other two plotted parameters are both not captured well. To
start of with, they both have a p-value of 0, which means that neither
of them accurately capture the observed data. The geodesic distribution
itself has values that are higher than they are supposed to be over all
distances, granting it a p-value of 0. This means that the simulated
network is more connected than the real one, and that the real one has
more of a community type structured one, meaning there are more separate
clusters that have only few edges going out / coming in. The Triad
consensus also has a p-value of 0, and as seen in the graph, while it
correctly predicts some kinds of triads, it is completely wrong by
overestimating some and underestimating others. It makes it clear that
some important terms for the development of the network are still left
out and not being considered.

\subsubsection{3.2.4}\label{section-19}

Hypothesis iv: Popular students tend to increase or maintain their level
of alcohol consumption: This hypothesis is not supported by the data.
The estimate for the effect of indegree on alcoholbeh is 0.0703, which
is positive but statistically insignificant given the standard error of
0.0645. Because the estimate is not at least twice the standard error,
we can't confidently conclude that popularity in the friendship network
drives alcohol consumption. While the direction suggest popular students
might consume more, the effect is not statistically distinguishable from
zero. Hypothesis v: Students adjust their alcohol level to match their
friends: This hypothesis is strongly supported, as the alcoholbeh
average similarity effect has a high positive coefficient of 6.9119 and
is statistically significant, as it has a standard error of 1.8432. This
indicates that students change their drinking habits to minimize the
distance between themselves and the average consumption level of their
friends. The high behavior distribution with a p-value of 0.502 confirms
that the model accurately captures these behavioral shifts.

Revisiting hypotheses i-iii: Hypothesis i: Students tend to be friends
with popular pupils: In the co-evolution model, the indegree-popularity
effect is now negative with a value of -0.2400 and significant as its
standard error is 0.0761. This differs significantly from the finding in
(1), where the effect was nearly zero and insignificant. This indicates
that when we are accounting for alcohol co-evolution, students actually
show a slight preference against forming ties with the most popular
peers. Hypothesis ii: Students tend to be friends with pupils with
similar alcohol consumption to theirs: The support for this hypothesis
remains strong, but the magnitude has increased significantly. The
alcoholbeh similarity coefficient rose from 0.494 in (1) to 1.1924 now.
This suggests that once we account for the fact that friends influence
each others drinking habits, the tendency to choose friends based on
pre-existing alcohol similarity appears to be even stronger than
previously thought. Hypothesis iii: Students tend to be friends with
students who live in the same neighborhood (living nearby): For this
hypothesis the findings remain consistent and significant. The
logdistance coefficient is -0.1950 with a standard error of 0.0461,
which closely matches the -0.185 we found in (1). The negative value
confirms that as the distance between homes increases, the likelihood of
friendship decreases, which supports our hypothesis that students from
the same neighborhoods tend to be friends. However the geodesic
distribution still has a p-value of 0, which tells us that while the
distance of students is important, it is not sufficient to fully explain
the network's overall community structure.

\subsubsection{3.2.5}\label{section-20}

We have strong evidence for both selection and influence processes
operating simultaneously within the co-evolution of the friendship
network and alcohol consumption. Selection processes can be seen by the
alcoholbeh similarity in the Network Dynamics part of the model, which
yields a positive estimate of 1.1924 with a standard error of 0.3403.
Since this coefficient is more than three times its standard error, it
is statistically significant, which indicates that students actively
select friends who have similar drinking habits to their own. This
selection based on behavioral homophily is a key driver of the
friendship network's structural evolution. Next to the selection
processes the model also provides strong evidence for influence
processes in the Behavior Dynamics part of the model. Here it shows a
large estimate of 6.9119, that is stastically significant with a
standard error of 1.8432, which confirms that students do not just
choose similar friends, but also adapt their own alcohol consumption
over time to align with the average behavior of their existing friend
group (as also discussed in Task 3.2.4).

\subsection{3.3}\label{section-21}

There are more nodes with a low geodesic distance in the simulated model
than in the real observed one, meaning that the simulated one does not
create as strongly isolated clusters as the observed model has. A
simulated model could achieve better scores here, if it encourages
isolated clusters more. This means we would require that the current
transitivity effect becomes stronger. At the same time, we need the
students to be more likely to cut a tie when they do not share many
common friends and are not connected well to the node otherwise. This
should lead to stronger community forming. To futher get all the effects
on transitivity, we should add the \(transRecTrip\) and the \(gwesp\)
terms.

The indegree distribution already closely follows the real world model,
so it does not require any changes to it, since it currently captures
the model very well.

The outdegree distribution could be more accurate. It has too many
students with a low outdegree, and too few students with a high
outdegree. This too can be due to a too weak community clustering,
although that should have a rather small effect on it, the model might
just have an overall too high penalty for outgoing ties which would need
to be reduced. We might need to include an additional parameter like
\(outAct\) to capture more details on the network structure.

\subsection{3.4}\label{section-22}

An additional hypothesis to further look into the co-evolution of social
structures and alcohol consumption could be that the relationship
between alcohol consumption and popularity is moderated by gender,
suggesting that a high alcohol consumption increases the popularity of
males but decreases it for females. To operationalize the corresponding
effects in the evaluation function, one would collect gender as a binary
attribute and specify a behavior-related popularity effect interacted
with a gender-ego effect. This tests if the ``social reward'' of
drinking depends on the gender and would not treat all students as a
homogeneous group, as there are now genders associated with each
student. A second additional hypothesis to be tested would be that the
influence of peers on alcohol consumption is influences by the frequency
of their physical interactions outside of school. More specifically this
means that students are more likely to assimilate to the alcohol levels
of friends with whom they have a lot of physical interactions. This
could be operationalized by collecting the contact frequency of friends
as a dyadic covariate and implementing a weighted average similarity
effect in the behavior dynamics part. Rather than the broad alcoholbeh
average similarity coefficient, this would weight the influence of each
friend by the reported frequency of their interactions.

\end{document}
